% Note du modèle mathématique unifié de l'idée intéressante et de sa naissance (v1)
% Compilation : xelatex deux fois.
\documentclass[11pt,a4paper]{article}

\usepackage{fontspec}
\usepackage{polyglossia}
\setdefaultlanguage{french}
\setmainfont{Linux Libertine O}
\setsansfont{DejaVu Sans}
\setmonofont[Scale=0.85]{DejaVu Sans Mono}

\usepackage{amsmath,amssymb,amsthm,mathtools}
\usepackage[a4paper,margin=2.5cm]{geometry}
\usepackage{microtype}
\usepackage{csquotes}
\usepackage{booktabs,tabularx,longtable,array}
\usepackage{enumitem}
\usepackage[most]{tcolorbox}
\usepackage{tikz}
\usetikzlibrary{arrows.meta,positioning,calc,shapes.geometric,decorations.pathmorphing}
\usepackage{caption}
\usepackage{xurl}
\usepackage[hidelinks]{hyperref}

\setlist{itemsep=2pt,topsep=4pt}
\renewcommand{\arraystretch}{1.18}
\setlength{\emergencystretch}{2em}

\theoremstyle{definition}
\newtheorem{definition}{Définition}[section]
\newtheorem{exemple}[definition]{Exemple}
\newtheorem{remarque}[definition]{Remarque}
\newtheorem{test}[definition]{Épreuve}
\theoremstyle{plain}
\newtheorem{proposition}[definition]{Proposition}
\newtheorem{corollaire}[definition]{Corollaire}
\renewcommand{\qedsymbol}{$\square$}

\newcommand{\R}{\mathcal{R}}
\newcommand{\Inc}{\mathcal{I}}
\newcommand{\dom}{\operatorname{dom}}
\newcommand{\chemin}[1]{\nolinkurl{#1}}
\newcommand{\vg}{\,{,}\,}

% Encadré des trois statuts
\newtcolorbox{statuts}[1][]{enhanced,breakable,colback=black!3,colframe=black!55,
  boxrule=0.6pt,arc=1.5pt,left=6pt,right=6pt,top=4pt,bottom=4pt,
  fonttitle=\bfseries\small,title={#1}}
\newcommand{\demontre}{\par\noindent\textbf{Ce qui est démontré.}~}
\newcommand{\mesure}{\par\smallskip\noindent\textbf{Ce qui est mesuré.}~}
\newcommand{\conjecture}{\par\smallskip\noindent\textbf{Ce qui est conjecturé.}~}

\newtcolorbox{attente}{enhanced,breakable,colback=orange!6,colframe=orange!60!black,
  boxrule=0.9pt,arc=2pt,left=7pt,right=7pt,top=5pt,bottom=5pt,
  fonttitle=\bfseries,title={En attente de Fable, après le 28 novembre 2026 à {\NoAutoSpacing 04:00} UTC}}

\hypersetup{pdftitle={Modèle mathématique unifié de l'idée intéressante et de sa naissance (v1)},
  pdfauthor={Claude Opus 5.5, sous-agent de la session des questions du directeur},
  pdflang={fr}}

\title{\textbf{Un modèle mathématique unifié de l'idée intéressante\\ et de sa naissance}\\[4pt]
\large Note de travail, version 1}
\author{Claude Opus 5.5 (effort élevé), sous-agent de la session fille des questions du directeur}
\date{24 novembre 2026, {\NoAutoSpacing 16:52} UTC}

\begin{document}
\maketitle

\begin{center}
\begin{minipage}{0.92\textwidth}\small
\textbf{Nature.} Pièce de travail (classe C), livrable (b) du dossier
\chemin{docs/notes/questions-directeur-2026-11-24-v1/}. Aucune dépense, aucun appel à un modèle
externe, aucune course. Rien n'a été ouvert sous \chemin{coffre/}, \chemin{docs/verdicts/},
\chemin{runs/}, \chemin{diag/}, \chemin{donnees/}.

\textbf{Matière.} Les réponses aux questions 1, 2, 3, 4, 5, 7 et 9 du directeur et leurs errata,
qui priment ; la contre-lecture ; les trois compléments transmis par la mère ; la réponse de Fable
(partie 2) et les réponses de Carbon. Empreintes en annexe~\ref{annexe:sources}. Tout chiffre de
la note renvoie à une pièce empreintée, à une page d'une source, ou à un calcul refait ici et
dit tel ; sinon il est marqué \enquote{à vérifier}.

\textbf{Règle de lecture.} Rien n'est présenté comme neuf (section~\ref{sec:neuf}). Le mot
\enquote{démontré} ne qualifie qu'un énoncé dont les hypothèses sont écrites et la preuve donnée ou
citée ; presque toutes les preuves de la note sont courtes et élémentaires, et leur intérêt est de
fixer les hypothèses, non d'établir un résultat profond.
\end{minipage}
\end{center}

\begin{attente}
Trois questions du directeur n'ont pas encore de réponse : l'agent Fable qui les traitait a été
arrêté par la limite hebdomadaire de son quota (faux échec, panne technique ; travail partiel
versé dans \chemin{partiel-fable-q6-q8-coupe-quota-v1/}, non relu, non réemployé ici).
\begin{itemize}
\item \textbf{Question 6} : peut-on \emph{démontrer}, à partir du modèle unifié, comment maximiser la
probabilité d'une découverte ; conditions d'émergence ; critères qui aident l'organisme
d'idéation à faire de vraies découvertes profondes.
\item \textbf{Question 8} : la \enquote{plasticité} d'une région du graphe d'un domaine (sa capacité de
restructuration compatible avec les relations existantes) prédit-elle une zone de découverte,
et quel est son lien avec les anomalies, les équilibres et les bifurcations ?
\item \textbf{Question 10} : que peut-on prédire du rôle de l'objet manquant à partir de la
\emph{forme} de l'anomalie, démonstration à l'appui ?
\end{itemize}
\textbf{Ce que la note en dit seulement : le cadre, sans démonstration neuve.}
Pour la question 6, la note pose les objets et le taux de naissance sur lesquels une politique de
recherche serait optimisée (sections~\ref{sec:objets} à \ref{sec:cinetique}) ; les théorèmes
classiques que la demande reformulée nomme (indice de Gittins, règle de Weitzman, théorie de la
recherche optimale, garantie gloutonne des fonctions sous-modulaires, criticalité d'un processus
de branchement) ne sont ni énoncés ni vérifiés ici. Pour la question 8, la note donne seulement les
deux perturbations (des faits, de la correspondance) qui serviraient de briques à une définition de
la plasticité (section~\ref{sec:perturbations}) ; les trois définitions candidates de la demande
reformulée (dégénérescence du paysage des partitions, susceptibilité de la structure optimale à
l'ajout d'une relation, mollesse d'une énergie de frustration) ne sont pas traitées. Pour la
question 10, la note donne seulement le \enquote{profil} qu'une correspondance impose à un objet
manquant (définition~\ref{def:profil}) : c'est le type des relations que l'objet doit porter, non le
type de l'objet ; la thèse \enquote{la forme de l'anomalie prédit le type de l'objet} reste
entière. Les trois réponses de Fable primeront sur tout ce que la note dit à leur sujet.
\end{attente}

\tableofcontents

%======================================================================
\section{Guide de lecture}
\label{sec:guide}

\subsection{Trois questions qu'il ne faut pas confondre}
La réponse de Fable \cite{fable} sépare trois questions, et toute la note les garde séparées.
\begin{enumerate}[label=(\alph*)]
\item \textbf{Intéressant pour un public} : une propriété relative à un état de croyances daté ;
une thèse intéresse quand elle nie une part de ce que son public tient pour acquis, ennuie quand
elle le confirme, paraît absurde quand elle en nie tout le fond \cite{davis1971}.
\item \textbf{Vrai et fécond} : une propriété du monde, connue après coup.
\item \textbf{Produire} : un processus de recherche, avec son taux.
\end{enumerate}
Elles divergent : des évaluateurs affectés au hasard notent plus bas les projets très
nouveaux \cite{boudreau2016}, alors que la surprise prédit l'impact ultérieur
\cite{shievans2023}. Dans la note, (a) est l'intérêt $I$ estimé avant (section~\ref{sec:surprise}),
(b) la portée réalisée et la fécondité (sections~\ref{sec:surprise} et \ref{sec:fecondite}), (c) le
taux de naissance $\lambda$ (sections~\ref{sec:objets} et \ref{sec:cinetique}).

\subsection{Trois statuts}
Chaque section finit par un encadré à trois rubriques. \emph{Démontré} : énoncé à hypothèses
écrites, preuve donnée ou citée. \emph{Mesuré} : chiffre sorti d'une pièce du dossier, à sa page
ou à son empreinte. \emph{Conjecturé} : énoncé plausible, avec l'épreuve qui le réfuterait. Une
épreuve décrite ici n'est ni préenregistrée ni lancée : toute épreuve décisive exige un
préenregistrement scellé et contre-lu, et un accord sur devis.

\subsection{Conventions de notation}
L'erratum transversal \cite{erratumkappa} fixe une notation pour tout le dossier ; la note la suit
(tableau~\ref{tab:notations}). Deux choix propres à la note : le type d'une relation s'écrit $r$
(et non $\rho$, homonyme chez Fable), et le coefficient de la surprise dans la crédibilité s'écrit
$\vartheta$ (Fable l'écrit $\rho$).

\begin{table}[ht]
\centering\small
\caption{Notations du dossier, d'après l'erratum transversal \cite{erratumkappa}.}
\label{tab:notations}
\begin{tabularx}{\textwidth}{@{}lX@{}}
\toprule
Symbole & Sens \\
\midrule
$c(\varphi)$ & couverture : part des relations du voisinage source dont l'image existe avec le même type \\
$\kappa(\varphi)$ & \textbf{part des conflits} : part des relations dont l'image existe avec un type incompatible ; jamais une \enquote{cohérence} \\
$K_r$, $Z_r$ & score serrure-clé $\max_\varphi[c-\kappa]$ au rayon $r$, et son écart \emph{réduit} sous recâblage \\
$c_C$, $\chi_C$, $\gamma_C$, $K_C$ & couverture pondérée, taux de conflits, pénalité et score serrure-clé \emph{chez Carbon} \\
$\psi$, $k_0$, $w_m$ & coefficient de la voie catalytique, constante de rencontre, poids de la méthode $m$ (question 3) \\
$\beta_E$, $\delta_E$ & coefficients de la barrière $E_{ab}=E_0-\beta_E Z_{ab}-\delta_E P_{ab}$ \\
$G(t)$, $\Delta G(t)$ & part des auteurs dans la plus grande composante, son excès sur le nul (question 4) \\
$P_{ab}$, $\delta$ & terme des problèmes ouverts et son coefficient dans le taux de naissance \\
$r$ & type d'une relation $(x,r,y)$ \\
$\vartheta$ & coefficient de la surprise $S$ dans la crédibilité $q$ \\
\bottomrule
\end{tabularx}
\end{table}
L'accord entre lecteurs (kappa pondéré, alpha de Krippendorff) s'écrit toujours en toutes lettres.

%======================================================================
\section{Les objets : l'hypergraphe fusionné}
\label{sec:objets}

L'objet commun du dossier est \emph{l'hypergraphe fusionné} : celui de la question~7, section~5
\cite{q7}, qui réunit le cœur de Fable \cite{fable}, la résolvante et la couche typée de
Carbon \cite{carbonmeta,carbonhomo}, plus des poids, des anomalies empiriques, un ancrage et un
couplage au taux de naissance. La comparaison qui y conduit (sept critères fixés d'avance) conclut
qu'aucun des deux modèles ne domine \cite[section~4]{q7}, ce que l'erratum transversal fixe
pour tout le dossier \cite[section~2]{erratumkappa}. Carbon n'écrit nulle part
\enquote{hypergraphe} ; le mot est de Fable et de Neon \cite[section~2]{q7}.

\begin{definition}[Hypergraphe orienté]
\label{def:hypergraphe}
Un \emph{hypergraphe orienté} sur un ensemble de nœuds $V$ est un ensemble d'\emph{hyperarcs}
$e=(T(e),h(e))$, où la queue $T(e)\subseteq V$ est un ensemble fini non vide de prémisses et la tête
$h(e)\in V$ une conclusion \cite{gallo1993}. Un arc ordinaire est un hyperarc à queue singleton.
Un \emph{hyperchemin} de $X$ vers $y$ est une suite d'hyperarcs dont chaque queue est contenue dans
$X$ augmenté des têtes précédentes, la dernière tête étant $y$.
\end{definition}

\begin{definition}[Relations typées, alphabet et incompatibilité]
\label{def:alphabet}
Un \emph{alphabet des relations} est un ensemble fini $\R$ de types (\enquote{est loi sur},
\enquote{active}, \enquote{inhibe}, \enquote{divise}\dots), muni d'une relation symétrique
d'\emph{incompatibilité} $\Inc\subseteq\R\times\R$ (par exemple \enquote{active} et
\enquote{inhibe}, \enquote{vérifie} et \enquote{réfute}). Une \emph{structure relationnelle typée}
est un couple $G=(V,E)$ avec $E\subseteq V\times\R\times V$ ; chaque $(x,r,y)\in E$ est une
\emph{relation} de source $x$, de but $y$, de type $r$. Les types sont communs à tous les
domaines \cite[partie~2]{fable} ; seuls les noms des nœuds sont propres à un domaine.
\end{definition}

\begin{definition}[Hypergraphe fusionné daté]
\label{def:fusionne}
À chaque date $t$, l'\emph{hypergraphe fusionné} $\mathcal{H}_t$ porte, sur les concepts $V_t$
(dont les cadres $M\subseteq V_t$ : méthodes, instruments, formalismes), quatre couches :
\begin{enumerate}
\item \textbf{usage} $U_t$ : chaque document antérieur à $t$ est un hyperarc qui relie l'ensemble de
ses concepts (cooccurrence) ;
\item \textbf{structure} $R_t$ : relations typées, signées, orientées (définition~\ref{def:alphabet}),
les équations devenant des arbres d'opérateurs \cite[partie~2]{fable} ;
\item \textbf{humains} $A_t$ : qui a travaillé sur quoi (matrice d'incidence auteurs-concepts) ;
\item \textbf{problèmes} $Q_t$ : questions ouvertes et anomalies, hyperarcs \enquote{prémisses
$\to$ cible} pondérés (définitions~\ref{def:problemes} et \ref{def:poids}).
\end{enumerate}
Toute grandeur indexée par $t$ se calcule sur les seules données antérieures à $t$.
\end{definition}

\begin{definition}[Couche des problèmes]
\label{def:problemes}
D'après la question~7, section~5 \cite{q7} :
\begin{itemize}
\item un \emph{résultat} est un hyperarc établi ; une \emph{question} est un hyperarc demandé, sans
témoin connu, noté $(\Gamma,\tau,\text{date},\text{auteurs},\text{statut})$ (le séquent
\enquote{$\Gamma\vdash^{?}\tau$} de Carbon \cite[section~3]{carbonmeta}) ; une \emph{réduction} est un
hyperarc entre problèmes ; une \emph{dérivation} est un hyperchemin ;
\item une \emph{anomalie} oppose un hyperarc de prédiction (théorie $T$, hypothèses auxiliaires
$H\to$ grandeur $q$) à un hyperarc de mesure (instrument $M\to q$) ; son écart est
\[ z=\frac{|\text{prédit}-\text{mesuré}|}{\sqrt{\sigma_p^2+\sigma_m^2}}, \]
avec ouverture à $z\ge 3$ (seuil à geler), une durée et un mode de fermeture (idée, remesure,
recalcul) ;
\item un \emph{objet en gestation} est, en mathématiques, une macro des sources définie par plusieurs
groupes non cosignataires, à même arbre d'expansion près, avec ses axiomes partiels.
\end{itemize}
\end{definition}

\begin{definition}[Poids d'un problème]
\label{def:poids}
Le poids $w_j(t)$ du problème $j$ combine cinq facteurs mesurés avant $t$ \cite[section~5]{q7} :
(1) l'écart $z_j$ ; (2) la résistance : groupes d'auteurs disjoints qui citent $j$ sans le fermer ;
(3) la portée : problèmes que la solution de $j$ ferait avancer par les réductions ; (4) la
centralité des prémisses mises en cause, rapportée avec et sans (elle ramène le degré) ; (5) la
maturation : croissance des cadres employés par les documents ancrés à $j$. Règle fixée d'avance :
$\log w_j$ est la moyenne des rangs centrés réduits des cinq facteurs, par champ et par année ;
fermeture : $w_j=0$, mode consigné.
\end{definition}

Les facteurs du poids sont occupés depuis Laudan 1977 (discordance, âge, généralité ; pages
rapportées par Moser \cite{moser1991}, livre non vu \cite{erratamineurs}) et Kuhn
\cite[p.~52 et 81-82]{kuhn1962} ; restent leur mesure datée et le couplage.

\begin{definition}[Taux de naissance fusionné]
\label{def:taux}
Pour une paire de concepts $(a,b)$ jamais reliée avant $t$, le taux instantané de naissance du lien
est
\begin{equation}\label{eq:taux}
\lambda_{ab}(t)=p^0_t(a,b)\,\exp\bigl(\beta Z_{ab}+\gamma F_t(a,b)+\delta P_t(a,b)\bigr),
\end{equation}
où $p^0_t(a,b)$ est la probabilité de surface (modèle appris avant $t$ sur la proximité d'usage, la
popularité et les experts communs), $Z_{ab}$ le score serrure-clé réduit
(section~\ref{sec:correspondance}), $F_t(a,b)=\sum_{k\in M}\Delta h_k(t)\,s_a(k)\,s_b(k)$ le terme
des cadres ($\Delta h_k$ : croissance récente de la part des documents qui emploient le cadre $k$ ;
$s_a(k)$ : affinité de $a$ pour $k$ mesurée avant $t$) \cite[partie~2]{fable}, et
\begin{equation}\label{eq:P}
P_t(a,b)=\log\Bigl(1+\sum_j w_j(t)\,\mathbf{1}\bigl[\text{un parent ancré à }j\text{, l'autre non}\bigr]\Bigr)
\end{equation}
le terme des problèmes \cite[section~5]{q7}. La forme $\log\lambda=\log p^0+\beta Z+\gamma F+\delta P$
concorde dans les questions 3, 5, 7 et 9 \cite[contradictions, point~3]{contrelecture}.
\end{definition}

\begin{remarque}[Lecture des coefficients]
L'hypothèse \enquote{serrure et clé} dit $\beta>0$, l'hypothèse du champ $\gamma>0$, la couche des
problèmes $\delta>0$ à surface, structure et cadres égaux. La question~3 montre que $F$ se lit mieux
comme un catalyseur que comme un champ au sens physique (section~\ref{sec:cinetique}) : le nom
\enquote{champ} est gardé par usage. Un déplacement uniforme de toutes les barrières (effet
d'époque) est le taux de base $\lambda_0(t)$ d'un modèle de Cox, déjà absorbé.
\end{remarque}

\begin{exemple}[Une anomalie fermée par recalcul]
Le moment magnétique du muon était à 4,2 écarts types du modèle standard en 2021 \cite{abi2021}
(résumé) ; une mise à jour du calcul théorique ramène l'écart à $38(63)\times10^{-11}$, \enquote{no
tension}, en 2025 \cite{aliberti2025} (résumé). Dans la couche des problèmes : un hyperarc de
prédiction contre un hyperarc de mesure, ouvert à $z\ge3$, fermé par recalcul. La question~7 en
tire une prédiction (P4) : une idée née sur une anomalie fermée ensuite par recalcul ou remesure
est moins féconde \cite[section~5]{q7}.
\end{exemple}

\begin{figure}[ht]
\centering
\begin{tikzpicture}[x=1cm,y=1cm,
  nd/.style={circle,fill=black,inner sep=1.3pt},
  lab/.style={font=\small\itshape,anchor=west}]
\foreach \yy/\nm in {0/humains $A_t$,1.6/usage $U_t$,3.2/structure $R_t$,4.8/problèmes $Q_t$}{
  \draw[black!45,fill=black!4] (0,\yy) -- (7,\yy) -- (8,\yy+0.9) -- (1,\yy+0.9) -- cycle;
  \node[lab] at (8.2,\yy+0.45) {\nm};}
% structure
\node[nd] (s1) at (2,3.6) {}; \node[nd] (s2) at (3.5,3.75) {}; \node[nd] (s3) at (5,3.55) {}; \node[nd] (s4) at (6.2,3.8) {};
\draw[-{Stealth[length=4pt]}] (s1) -- node[above,font=\tiny]{$r$} (s2);
\draw[-{Stealth[length=4pt]}] (s2) -- node[above,font=\tiny]{$r'$} (s3);
\draw[-{Stealth[length=4pt]},dashed] (s3) -- (s4);
% usage : hyperarête
\draw[black!60,fill=black!12,rounded corners=6pt] (1.7,1.85) -- (3.8,2.2) -- (3.2,2.4) -- (1.6,2.2) -- cycle;
\node[nd] at (2,2.05) {}; \node[nd] at (3.3,2.2) {}; \node[nd] at (2.6,2.25) {};
% humains
\node[nd] (h1) at (2.4,0.45) {}; \node[nd] (h2) at (4.6,0.5) {};
\draw (h1) -- (2,1.6); \draw (h2) -- (4.4,1.6);
% problème : hyperarc prémisses -> cible
\node[nd] (p1) at (2.2,5.2) {}; \node[nd] (p2) at (3.0,5.4) {}; \node[draw,circle,inner sep=1.6pt] (pc) at (5.6,5.3) {?};
\draw[thick] (p1) -- (4,5.3); \draw[thick] (p2) -- (4,5.3);
\draw[thick,-{Stealth[length=5pt]}] (4,5.3) -- node[above,font=\tiny]{question} (pc);
\draw[dotted,thick] (pc) -- (5,3.55);
\draw[dotted,thick] (p1) -- (2,3.6);
\end{tikzpicture}
\caption{Les quatre couches de l'hypergraphe fusionné. En haut, une question : hyperarc demandé,
sans témoin, ancré (pointillés) aux concepts de la couche de structure. Le couplage $\delta P$ fait
peser le problème sur la naissance des liens dont un seul parent lui est ancré.}
\label{fig:couches}
\end{figure}

\begin{statuts}[Section 2 : statuts]
\demontre Rien : la section définit.
\mesure Rien ici ; les chiffres cités (muon) sont lus sur résumé par la question~7.
\conjecture $\delta>0$ (prédiction P1 de la question~7) ; la rupture datée P2 (les naissances
ancrées à $j$ montent après son ouverture, retombent après sa fermeture) ; P3 à P5 \cite[section~5]{q7}.
Épreuve proposée : greffée sur l'expérience de Fable (PubTator~3, gel fin 2010, naissances
2011-2015), rapport de risques par écart type de $P$ au moins 1,10, borne basse à 95~\% au-dessus
de 1, effet nul sous permutation des signes ; réfutation si la borne haute reste sous 1,05, gardes
tenues \cite[section~6]{q7}. Garde de mesurabilité d'abord : sur la base de prédications SemMedDB,
58 des 2\,236 paires candidates étaient des contradictions apparentes, surtout expliquées par le
contexte \cite{rosemblat2019} (résumé).
\end{statuts}

%======================================================================
\section{La correspondance comme homomorphisme partiel}
\label{sec:correspondance}

\subsection{Définitions}

\begin{definition}[Voisinage typé]
Pour un concept $a$ et un rayon $r\in\{1,2\}$, le voisinage typé $R^r_a=(V^r_a,E^r_a)$ est la
sous-structure de $R_t$ formée des relations à distance au plus $r$ de $a$, \emph{noms des nœuds
effacés, types gardés} \cite[partie~2]{fable}. On note $N=|E^r_a|$.
\end{definition}

\begin{definition}[Correspondance]
Une \emph{correspondance} de $a$ vers $b$ est une injection partielle
$\varphi:V^r_a\rightharpoonup V^r_b$ de domaine $\dom\varphi$, avec $\varphi(a)=b$.
\end{definition}

\begin{definition}[Homomorphisme partiel ; transport des relations]
\label{def:homo}
Soit $D\subseteq E^r_a$. On dit que $\varphi$ est un \emph{homomorphisme partiel sur $D$} si, pour toute
$(x,r,y)\in D$, on a $x,y\in\dom\varphi$ et $(\varphi x,r,\varphi y)\in E^r_b$. L'application de
transport $\varphi_E:D\to E^r_b$, $(x,r,y)\mapsto(\varphi x,r,\varphi y)$, vérifie alors
\[ \mathrm{src}_B\circ\varphi_E=\varphi\circ\mathrm{src}_A,\qquad
   \mathrm{but}_B\circ\varphi_E=\varphi\circ\mathrm{but}_A,\qquad
   \mathrm{type}_B\circ\varphi_E=\mathrm{type}_A \quad\text{sur } D : \]
\emph{les relations transportées commutent avec $\varphi$}. C'est la seule \enquote{conservation}
que le dossier retient : le théorème de Noether exige une action et un groupe continu de
symétries, que l'espace des concepts n'a pas \cite[partie~2]{fable} ; Carbon arrive à la même
séparation \cite[section~3]{carbonhomo}.
\end{definition}

\begin{definition}[Les quatre cas d'une relation]
\label{def:quatrecas}
Pour $e=(x,r,y)\in E^r_a$ et une correspondance $\varphi$, exactement un des cas suivants a lieu
\cite[correction (2)]{complement89} :
\begin{enumerate}[label=(C\arabic*)]
\item \textbf{confirmée} : $x,y\in\dom\varphi$ et $(\varphi x,r,\varphi y)\in E^r_b$ ; elle compte dans $c$ ;
\item \textbf{contredite} : $x,y\in\dom\varphi$, non confirmée, et il existe $r'$ avec
$(\varphi x,r',\varphi y)\in E^r_b$ et $(r,r')\in\Inc$ ; elle compte dans $\kappa$ ;
\item \textbf{absente} : $x,y\in\dom\varphi$, ni confirmée ni contredite ; c'est une
\emph{prédiction de relation} $(\varphi x,r,\varphi y)$ ;
\item \textbf{extrémité sans image} : $x\notin\dom\varphi$ ou $y\notin\dom\varphi$ ; c'est la
\emph{prédiction d'un objet manquant}, un \enquote{trou}, comme les cases vides du tableau de
Mendeleïev.
\end{enumerate}
\end{definition}

Le cas (C3) contient aussi une relation d'un autre type, compatible, entre $\varphi x$ et
$\varphi y$ ; on peut l'en séparer si l'alphabet le demande, sans rien changer à ce qui suit.

\begin{proposition}[Partition]
\label{prop:partition}
Soient $c,\kappa,a,m$ les parts des relations de $E^r_a$ dans les cas (C1) à (C4). Alors
$c+\kappa+a+m=1$, et la même identité vaut pour des parts pondérées par des poids positifs
quelconques.
\end{proposition}
\begin{proof}
Les quatre conditions sont exclusives par construction (C4 se lit sur le domaine ; C1, C2, C3 se
lisent, pour une relation dont les deux extrémités ont une image, dans cet ordre, chacune excluant
les précédentes) et leur réunion épuise les cas. Toute mesure positive additive sur $E^r_a$,
normalisée, donne donc des parts de somme 1.
\end{proof}

\begin{proposition}[Homomorphisme exact]
\label{prop:homoexact}
L'ensemble $D_\varphi$ des relations confirmées est le plus grand $D$ sur lequel $\varphi$ est un
homomorphisme partiel. En particulier $c(\varphi)=1$ si et seulement si $\varphi$ est un
homomorphisme partiel sur tout $E^r_a$, et alors $\kappa=a=m=0$.
\end{proposition}
\begin{proof}
Une relation appartient à un $D$ admissible si et seulement si elle vérifie (C1) ; la réunion des
$D$ admissibles est donc $D_\varphi$, lui-même admissible. Si $c=1$, toutes les relations sont
confirmées ; la proposition~\ref{prop:partition} donne $\kappa=a=m=0$. La réciproque est la
définition.
\end{proof}

\begin{figure}[ht]
\centering
\begin{tikzpicture}[x=1cm,y=1cm,
  nd/.style={circle,draw,fill=white,inner sep=1.6pt,font=\small},
  rel/.style={-{Stealth[length=4pt]},thick},
  phi/.style={-{Stealth[length=4pt]},dashed,black!55}]
\node[font=\bfseries] at (1.3,4.3) {domaine $A$};
\node[font=\bfseries] at (7.3,4.3) {domaine $B$};
\node[nd] (x) at (0,2) {$x$};
\node[nd] (y1) at (2.4,3.6) {$y_1$}; \node[nd] (y2) at (2.6,2.4) {$y_2$};
\node[nd] (y3) at (2.6,1.2) {$y_3$}; \node[nd] (y4) at (2.4,0) {$y_4$};
\draw[rel] (x) -- node[above,font=\tiny,sloped]{$r_1$} (y1);
\draw[rel] (x) -- node[above,font=\tiny,sloped]{$r_2$} (y2);
\draw[rel] (x) -- node[above,font=\tiny,sloped]{$r_3$} (y3);
\draw[rel] (x) -- node[below,font=\tiny,sloped]{$r_4$} (y4);
\node[nd] (X) at (6,2) {$\varphi x$};
\node[nd] (Y1) at (8.6,3.6) {$\varphi y_1$}; \node[nd] (Y2) at (8.8,2.4) {$\varphi y_2$};
\node[nd] (Y3) at (8.8,1.2) {$\varphi y_3$};
\node[draw,dashed,circle,inner sep=1.6pt,font=\small] (Y4) at (8.6,0) {?};
\draw[rel] (X) -- node[above,font=\tiny,sloped]{$r_1$} (Y1);
\draw[rel,red!70!black] (X) -- node[above,font=\tiny,sloped]{$r_2'$} (Y2);
\draw[phi] (x) to[bend left=12] (X);
\draw[phi] (y1) to[bend left=10] (Y1);
\draw[phi] (y2) to[bend left=6] (Y2);
\draw[phi] (y3) to[bend right=6] (Y3);
\node[font=\small,anchor=west] at (9.5,3.6) {(C1) confirmée : compte dans $c$};
\node[font=\small,anchor=west] at (9.5,2.4) {(C2) $(r_2,r_2')\in\Inc$ : dans $\kappa$};
\node[font=\small,anchor=west] at (9.5,1.2) {(C3) absente : prédiction de relation};
\node[font=\small,anchor=west] at (9.5,0) {(C4) sans image : objet manquant};
\end{tikzpicture}
\caption{Les quatre cas d'une relation $x\xrightarrow{r}y$ sous une correspondance $\varphi$
(flèches en tirets). En (C3), $\varphi$ prédit la relation $(\varphi x,r_3,\varphi y_3)$ ; en (C4),
elle prédit un objet de $B$ qui porterait la relation $r_4$ depuis $\varphi x$.}
\label{fig:quatrecas}
\end{figure}

\subsection{L'objet manquant et son profil}

\begin{definition}[Profil d'un objet manquant]
\label{def:profil}
Soit $y\in V^r_a\setminus\dom\varphi$. Son \emph{profil sous $\varphi$} est l'ensemble des relations
qu'un objet de $B$ devrait porter pour que les relations de $y$ soient confirmées :
\[ \Pi_\varphi(y)=\{(\varphi x,r,\star):(x,r,y)\in E^r_a,\ x\in\dom\varphi\}
\cup\{(\star,r,\varphi z):(y,r,z)\in E^r_a,\ z\in\dom\varphi\}, \]
où $\star$ désigne l'objet cherché.
\end{definition}

\begin{proposition}[Combler un trou]
\label{prop:combler}
Soit $y$ un nœud sans image dont toutes les relations ont leur autre extrémité dans $\dom\varphi$.
Ajoutons à $B$ un nœud neuf $y^\star$ portant exactement les relations de $\Pi_\varphi(y)$ (avec
$\star=y^\star$) et prolongeons $\varphi$ par $\varphi(y)=y^\star$. Alors la couverture croît de la
part des relations incidentes à $y$, et $\kappa$ ne change pas.
\end{proposition}
\begin{proof}
Seules changent de cas les relations incidentes à $y$ ; chacune passe de (C4) à (C1), puisque son
image est, par construction, une relation de même type de $B$. Aucune relation ne passe en (C2) :
les relations ajoutées ont toutes $y^\star$ pour extrémité et ne sont images que de relations
incidentes à $y$, toutes confirmées.
\end{proof}

\begin{remarque}[Le prix d'un trou]
La proposition~\ref{prop:combler} dit aussi le danger : ajouter des objets fait monter $c$ sans
rien coûter. Un score honnête doit donc faire payer l'objet ajouté, par exemple par sa longueur de
description \cite{rissanen1978} ; c'est l'objet de la \enquote{réification} de la question~9 et du
gain de longueur de description chez Carbon \cite[section~2.4]{carbon}. Le profil dit quelles
relations l'objet doit porter ; il ne dit pas ce qu'est l'objet (question~10, en attente).
\end{remarque}

\subsection{Le score serrure-clé, et le jouet de la question 9}

\begin{definition}[Score serrure-clé réduit]
$K_r(a,b)=\max_\varphi[c(\varphi)-\kappa(\varphi)]$, les relations entre relations pesant
davantage ; $Z_r=(K_r-\text{moyenne})/\text{écart-type}$ sous un recâblage qui préserve degrés et
types, par paire de domaines \cite[partie~2]{fable}. Critère : $Z_1$ et $Z_2$ élevés.
\end{definition}

\begin{exemple}[Permutations des racines et classes de formes quadratiques]
\label{ex:jouet}
La question~9 code les permutations des racines ($A$, huit relations typées) et les classes de
formes quadratiques binaires ($B$, six relations) \cite[section~3]{q9}. La correspondance
historique $\varphi$ a cinq relations confirmées, une contredite (la composition réfute la
commutativité dans $A$, la vérifie dans $B$), une absente (\enquote{le nombre de valeurs divise
l'ordre total} : prédiction de relation, juste) et une extrémité sans image (la résolvante) :
$c=5/8$, $\kappa=1/8$, $c-\kappa=0{,}500$. Sur les 362\,880 bijections des neuf nœuds appariables,
huit battent $\varphi$ (maximum $0{,}625$) en envoyant la commutativité sur un nœud sans relation
avec la composition : \emph{le conflit devient une absence, qui ne coûte rien}
\cite[section~5]{q9}. La conclusion tient sous la contrainte $\varphi(\text{composition})=
\text{composition}_B$ et pour les injections partielles \cite[mineur~1]{erratumq9}.
\end{exemple}

\begin{proposition}[Aucune pondération des cas ne fait de $\varphi$ l'unique optimum du jouet]
\label{prop:jouet}
Dans le jouet de l'exemple~\ref{ex:jouet}, soit un score $\Sigma(\varphi')$ qui ne dépend que des
nombres de relations dans les cas (C1) à (C4). Alors $\varphi$ n'est jamais l'unique maximum de
$\Sigma$. Pour le score $c-\kappa-\eta a$ avec $0\le\eta<1$, $\varphi$ est strictement battue.
\end{proposition}
\begin{proof}
La sortie de la version 2 du script \cite[section~2]{sortieq9v2} donne, pour la transposition qui
échange les images de l'associativité et de la commutativité (elle conserve le type des nœuds),
les mêmes nombres que $\varphi$ : cinq confirmées, une contredite, une absente, une sans image ; de
même pour l'échange du nombre de valeurs et de l'ordre total. Tout score fonction de ces nombres
leur donne la valeur de $\varphi$ : celle-ci n'est pas unique. Pour le score
$c-\kappa-\eta a$, la même sortie donne, pour les transpositions qui échangent la commutativité
avec le nombre de valeurs ou avec l'ordre total, cinq confirmées, aucune contredite, deux absentes,
une sans image ; leur score $(5-2\eta)/8$ dépasse celui de $\varphi$, $(4-\eta)/8$, si et seulement
si $\eta<1$.
\end{proof}

Leçon : faire payer l'absence ne suffit pas (à $\eta=1$, absence et conflit coûtent autant, et
l'on n'obtient qu'une égalité ; au-delà, une absence coûterait plus qu'un conflit) ; les types de nœuds non plus, puisque les deux transpositions
qui égalent $\varphi$ conservent les types. Dans ce jouet, la structure seule ne distingue pas
l'associativité de la commutativité \cite[section~5]{q9}. Il faut le nul du recâblage, des relations
plus nombreuses, ou un oracle extérieur (monde planté, assistant de preuve).

\subsection{Ce que Carbon propose}

Carbon écrit le même transport avec d'autres choix, que l'erratum transversal renomme pour ne pas
confondre son $\kappa$ (une couverture) avec celui de Fable (des conflits) \cite{erratumkappa}.
\begin{itemize}
\item \textbf{Couverture pondérée} \cite[section~2.2, l.~36-38]{carbon} :
\[ c_C(\varphi)=\frac{\sum_{(z_1,r,z_2)\in E_u(S)}\omega(z_1,z_2)\,\mathbf{1}\bigl[(\varphi z_1,\pi(r),\varphi z_2)\in\pi(E_v)\bigr]}{\sum_{(z_1,r,z_2)\in E_u}\omega(z_1,z_2)}, \]
où $S=\dom\varphi$, $E_u(S)$ les relations à deux
extrémités dans $S$, $\omega>0$ décroît avec la distance à $u$, et $\pi$ projette les types
disciplinaires sur un alphabet abstrait.
\item \textbf{Taux de conflits} : $\chi_C(\varphi)$, part des relations de $E_u(S)$ dont l'image porte
un type incompatible, rapportée à $|E_u(S)|$ et non à $|E_u|$.
\item \textbf{Score} \cite[l.~39-40]{carbon} : $K_C(u,v)=\max_\varphi c_C(\varphi)\,e^{-\gamma_C\chi_C(\varphi)}\in[0,1]$,
distance de structure $1-K_C$.
\item \textbf{Puits} \cite[l.~62]{carbon} : énergie d'alignement $E_C(\varphi)=-\log c_C(\varphi)+\gamma_C\chi_C(\varphi)$,
de sorte que $K_C=\max_\varphi e^{-E_C(\varphi)}$ ; la \enquote{profondeur} d'une analogie est
l'écart $\Delta E_C=E_C(\varphi_{(2)})-E_C(\varphi_{(1)})$ entre le second alignement et l'optimum.
\item \textbf{Fausse identité} : deux concepts confondus en surface se séparent quand
$|K_C(u,v)-K_C(u',v)|\gg0$ pour un concept sonde $v$.
\item \textbf{Critique de l'homomorphisme} \cite[verdict et section~1]{carbonhomo} : un homomorphisme
envoie une relation sur une relation, alors qu'une analogie envoie souvent une relation sur une
\emph{chaîne} ; d'où des foncteurs partiels entre catégories de chemins et, pour la détection, un
transport de Gromov et Wasserstein partiel calculé sur les résolvantes $(I-\alpha\tilde A)^{-1}$.
Les références de Carbon ne sont pas vérifiées ici.
\end{itemize}

\begin{proposition}[Quand les deux scores coïncident, et quand ils divergent]
\label{prop:fablecarbon}
Si $\omega\equiv1$, $\pi$ est l'identité et $\dom\varphi$ contient toutes les extrémités, alors
$c_C=c$ et $\chi_C=\kappa$. Même alors, les deux scores peuvent classer différemment deux
correspondances.
\end{proposition}
\begin{proof}
Sous ces hypothèses, $E_u(S)=E_u$ : les deux numérateurs et dénominateurs coïncident. Pour la
divergence, prenons $\gamma_C=1$, $\varphi_1$ avec $(c,\kappa)=(0{,}60\vg0{,}20)$ et $\varphi_2$
avec $(0{,}45\vg0)$. Fable : $c-\kappa$ vaut $0{,}40$ et $0{,}45$, $\varphi_2$ gagne. Carbon :
$c\,e^{-\kappa}$ vaut $0{,}60\,e^{-0{,}2}\approx0{,}491$ et $0{,}45$, $\varphi_1$ gagne (calcul
de l'auteur).
\end{proof}
Le choix de la forme (différence ou produit) n'est donc pas neutre ; il se fixe avant tout
préenregistrement. Les deux notions de \enquote{profondeur} diffèrent aussi : l'écart $\Delta E_C$
de Carbon mesure l'isolement du meilleur alignement, une grandeur de paysage, proche de la
largeur et de la rareté de la section~\ref{sec:perturbations} ; la profondeur du dossier est la
surprise retirée ailleurs (section~\ref{sec:surprise}).

\begin{proposition}[La résolvante voit les chaînes]
\label{prop:resolvante}
Soit $\tilde A\ge0$ une matrice d'adjacence pondérée et $0<\alpha<1/\varrho(\tilde A)$, où
$\varrho$ est le rayon spectral. Alors $(I-\alpha\tilde A)^{-1}=\sum_{k\ge0}\alpha^k\tilde A^k$, et
s'il existe un chemin $i\to k\to j$ de poids $\tilde A_{ik}\tilde A_{kj}>0$, l'entrée $(i,j)$ de la
résolvante est au moins $\alpha^2\tilde A_{ik}\tilde A_{kj}>0$, même sans relation directe $i\to j$.
\end{proposition}
\begin{proof}
Pour $\alpha\varrho(\tilde A)<1$ la série de Neumann converge et vaut l'inverse. Tous ses termes sont
positifs ou nuls ; le terme $k=2$ contient $\alpha^2\tilde A_{ik}\tilde A_{kj}$.
\end{proof}
C'est la seconde épreuve de la question~7 : dans un monde planté où une relation de $A$ correspond
à une chaîne de deux relations de $B$, le score $Z$ de Fable reste au nul, la résolvante retrouve
la correspondance \cite[section~4]{q7}.

\subsection{Deux perturbations distinctes}
\label{sec:perturbations}

\begin{definition}[Perturber les faits : robustesse]
Pour $\varepsilon\in(0,1)$, une \emph{perturbation des faits} retire chaque relation de $B$ (ou de $A$)
avec probabilité $\varepsilon$, ou en ajoute de fausses au même taux. La \emph{robustesse} de $\varphi$
est la loi de $c-\kappa$ sous ces perturbations, résumée par sa médiane et ses quantiles à 5~\% et
95~\%. Elle éprouve la \emph{profondeur} au sens de la tenue du score \cite[correction (3)]{complement89}.
\end{definition}

\begin{definition}[Perturber la correspondance : largeur locale et rareté au hasard]
\label{def:largeur}
Soit $d(\varphi,\varphi')$ le nombre minimal de transpositions qui mènent de $\varphi$ à $\varphi'$
(distance de Cayley). Pour une tolérance $\tau\ge0$ :
\[ L_\tau(\varphi;d)=\frac{|\{\varphi':d(\varphi,\varphi')=d,\ \Sigma(\varphi')\ge\Sigma(\varphi)-\tau\}|}{|\{\varphi':d(\varphi,\varphi')=d\}|},\qquad
   \mathrm{Rar}_\tau(\varphi)=\frac{|\{\varphi':\Sigma(\varphi')\ge\Sigma(\varphi)-\tau\}|}{|\{\varphi'\}|}. \]
$L_\tau(\varphi;1)$ et $L_\tau(\varphi;2)$ sont la \emph{largeur locale} (un puits \emph{étroit} a une
petite largeur locale) ; $\mathrm{Rar}_\tau$ est la \emph{rareté au hasard} \cite{erratumq9}.
\end{definition}

\begin{proposition}[Une transposition ne bouge le score que par les relations qu'elle touche]
\label{prop:lipschitz}
Soit $\Sigma=c-\kappa$ non pondéré, et $\varphi'$ obtenue en échangeant les images de deux nœuds
$u,v$. Soit $d_u$ le nombre de relations de $E^r_a$ incidentes à $u$. Alors
\[ |\Sigma(\varphi')-\Sigma(\varphi)|\le\frac{2(d_u+d_v)}{N}. \]
En conséquence, si $2(d_u+d_v)\le\tau N$ pour toute paire, alors $L_\tau(\varphi;1)=1$ : le puits
ne peut être localement étroit que si quelques nœuds portent une part notable des relations.
\end{proposition}
\begin{proof}
Le cas d'une relation $(x,r,y)$ ne dépend que de $\varphi x$ et $\varphi y$. L'échange ne modifie
ces images que si $x$ ou $y$ appartient à $\{u,v\}$ : au plus $d_u+d_v$ relations changent de cas.
Chacune contribue $+1$ (C1), $-1$ (C2) ou $0$ (C3, C4) au numérateur $N\Sigma$, donc change sa
contribution d'au plus 2. Pour des poids, remplacer les comptes par les poids.
\end{proof}

\begin{proposition}[Coût d'une recherche au hasard]
\label{prop:hasard}
Si l'on tire des correspondances uniformément et indépendamment, le nombre de tirages nécessaire
pour en trouver une à $\tau$ près de $\varphi$ avec probabilité $0{,}95$ est
$\lceil\ln0{,}05/\ln(1-\mathrm{Rar}_\tau)\rceil\approx3/\mathrm{Rar}_\tau$.
\end{proposition}
\begin{proof}
Le premier succès suit une loi géométrique de paramètre $\mathrm{Rar}_\tau$ :
\[ \Pr(\text{aucun succès en }n\text{ tirages})=(1-\mathrm{Rar}_\tau)^n, \]
et $\ln(1-x)\approx-x$ pour $x$ petit.
\end{proof}

\begin{table}[ht]
\centering\small
\caption{Largeur locale et rareté au hasard dans le jouet (tolérance $\tau=1/8$), d'après
l'erratum 1 de la question~9 \cite{erratumq9} et sa sortie \cite[section~1]{sortieq9v2}. Les deux
dernières lignes sont calculées ici par la proposition~\ref{prop:hasard}.}
\label{tab:largeur}
\begin{tabular}{@{}lrrr@{}}
\toprule
Distance à $\varphi$ & correspondances & à $\tau$ près de $\varphi$ & égalent ou battent $\varphi$ \\
\midrule
1 (une transposition) & 36 & 16 (0,444) & 7 (0,194) \\
2 & 546 & 99 (0,181) & 31 (0,057) \\
3 & 4\,536 & 272 (0,060) & 43 (0,009) \\
toutes & 362\,880 & 768 ($2{,}12\times10^{-3}$) & 100 ($2{,}76\times10^{-4}$) \\
\midrule
\multicolumn{4}{@{}l}{Tirages pour 95~\% de chances d'en trouver une à $\tau$ près : 1\,414 ; sous
$\varphi(\text{composition})$ fixée, 156.}\\
\bottomrule
\end{tabular}
\end{table}

Localement, le puits du jouet est \emph{large} ; ce qui est petit, c'est la part des bonnes
correspondances dans tout l'espace. Largeur locale et rareté au hasard sont deux grandeurs
distinctes, et la phrase \enquote{le puits est étroit} de la question~9 tombe \cite{erratumq9}. La
proposition~\ref{prop:lipschitz} explique en partie cette largeur : huit relations, plusieurs nœuds
de faible degré.

\begin{statuts}[Section 3 : statuts]
\demontre La partition en quatre cas et $c+\kappa+a+m=1$ (proposition~\ref{prop:partition}) ;
la caractérisation de l'homomorphisme exact (\ref{prop:homoexact}) ; combler un trou fait monter $c$
sans toucher $\kappa$ (\ref{prop:combler}) ; dans le jouet, aucun score fonction des seuls cas ne
rend $\varphi$ unique (\ref{prop:jouet}) ; les scores de Fable et de Carbon peuvent classer
différemment (\ref{prop:fablecarbon}) ; la résolvante voit les chaînes (\ref{prop:resolvante}) ;
la borne de la transposition (\ref{prop:lipschitz}) ; le coût d'une recherche au hasard
(\ref{prop:hasard}).
\mesure Sur le jouet : les comptes du tableau~\ref{tab:largeur} ; sous perturbation des faits
(retrait à 0,15, 2\,000 tirages), $c-\kappa$ vaut 0,430 en moyenne, entre 0,250 et 0,625 aux quantiles
5~\% et 95~\% : retirer un fait peut \emph{relever} le score en cachant un conflit
\cite[section~5]{q9}.
\conjecture Que, sur de vrais domaines, les correspondances qui ont mené à des découvertes soient
à la fois profondes et étroites (section~\ref{sec:puits}).
\end{statuts}

%======================================================================
\section{Coordonnées spectrales et applications fonctionnelles}
\label{sec:spectral}

\subsection{Les harmoniques d'un graphe}
Soit un graphe non orienté pondéré, de matrice d'adjacence symétrique $W\ge0$ et de matrice des
degrés $D=\operatorname{diag}(W\mathbf{1})$. Le \emph{laplacien} est $L=D-W$, et le laplacien
normalisé $\mathcal{L}=D^{-1/2}LD^{-1/2}$ \cite{chung1997}. Pour toute fonction $f$ sur les nœuds,
\[ f^\top Lf=\sum_{\{i,j\}}W_{ij}(f_i-f_j)^2\ \ge0. \]
$L$ est symétrique positif ; il s'écrit $L=\Phi\Lambda\Phi^\top$ avec $\Phi$ orthogonale (ses colonnes,
les \emph{modes propres}) et $\Lambda=\operatorname{diag}(\lambda_1\le\dots\le\lambda_n)$. Un mode de
petite valeur propre varie peu le long des arêtes : c'est une \enquote{basse fréquence}. La
multiplicité de la valeur propre 0 est le nombre de composantes connexes, et la deuxième valeur
propre mesure la connexité algébrique \cite{fiedler1973}. Une \emph{coordonnée spectrale} d'un nœud
est sa valeur dans un mode : une dimension est un \emph{mode de structure}, non une propriété d'un
nœud \cite[point (1)]{complementspectral}.

\subsection{Une correspondance comme matrice entre bases propres}

\begin{definition}[Application fonctionnelle]
\label{def:fmap}
Soient deux graphes $A$ et $B$ à $n$ nœuds, de bases propres orthonormées $\Phi_A,\Phi_B$, et une
bijection $\varphi:V_A\to V_B$ de matrice $P$ ($P_{xy}=1$ si $y=\varphi(x)$). L'application
$g\mapsto Pg=g\circ\varphi$ ramène les fonctions de $B$ sur $A$. Dans les bases propres elle s'écrit
\[ C=\Phi_A^\top P\,\Phi_B,\qquad C_{ij}=\langle\phi^A_i,\,\phi^B_j\circ\varphi\rangle . \]
C'est la représentation d'Ovsjanikov et collègues \cite{ovsjanikov2012} : les contraintes sur
$\varphi$ deviennent linéaires en $C$, l'estimation tient en \enquote{a single linear solve}
\cite[p.~1-2]{ovsjanikov2012} ; pour des formes presque isométriques, $C$ est \enquote{close to being
sparse and diagonally dominant}, sans être diagonale \cite[p.~3]{ovsjanikov2012} ; la commutation
avec des opérateurs conservés, $R_FC=CS_F$, donne des équations linéaires en $C$
\cite[p.~6]{ovsjanikov2012}. (Leur convention transporte les fonctions dans l'autre sens ; rien
ne change à ce qui suit.)
\end{definition}

\begin{proposition}[Isomorphisme : $C$ commute avec les spectres]
\label{prop:fmapiso}
Si $\varphi$ est un isomorphisme de graphes pondérés ($W^B_{\varphi x,\varphi y}=W^A_{xy}$), alors
$\Lambda_AC=C\Lambda_B$, donc $C_{ij}=0$ dès que $\lambda^A_i\ne\lambda^B_j$. Si de plus le spectre est
simple, $C$ est diagonale, à coefficients $\pm1$.
\end{proposition}
\begin{proof}
L'isomorphisme s'écrit $W_A=PW_BP^\top$ ; comme $P$ est une permutation, $D_A=PD_BP^\top$ et
$P^\top P=I$, d'où $L_AP=PL_B$. Alors
\[ \Lambda_AC=\Phi_A^\top L_AP\Phi_B=\Phi_A^\top PL_B\Phi_B=C\Lambda_B, \]
soit $(\lambda^A_i-\lambda^B_j)C_{ij}=0$. Deux graphes isomorphes ont le même spectre ; s'il est simple,
seuls les $C_{ii}$ peuvent être non nuls ; $C$ est orthogonale (produit de trois matrices
orthogonales), donc $C_{ii}=\pm1$.
\end{proof}

\begin{proposition}[Presque isomorphisme : $C$ presque diagonale, sauf entre modes proches]
\label{prop:fmapproche}
Soit $\varepsilon=\|L_AP-PL_B\|_2$ le défaut de commutation de $\varphi$. Pour tous $i,j$ tels que
$\lambda^A_i\ne\lambda^B_j$,
\[ |C_{ij}|\le\frac{\varepsilon}{|\lambda^A_i-\lambda^B_j|}. \]
\end{proposition}
\begin{proof}
$(\Lambda_AC-C\Lambda_B)_{ij}=(\lambda^A_i-\lambda^B_j)C_{ij}$, et cette entrée vaut
$(\phi^A_i)^\top(L_AP-PL_B)\phi^B_j$, majorée en valeur absolue par $\varepsilon$ puisque les deux
vecteurs sont unitaires.
\end{proof}

La proposition~\ref{prop:fmapproche} donne un sens précis à \enquote{presque diagonale si $\varphi$
conserve la structure} \cite[point (1)]{complementspectral} : la masse hors diagonale est bornée
par le défaut de commutation \emph{divisé par l'écart spectral}. Quand deux valeurs propres sont
proches, la borne ne dit plus rien.

\begin{exemple}[Le cycle : des modes égaux, une matrice par blocs]
\label{ex:cycle}
Sur le cycle à $n$ nœuds, les valeurs propres du laplacien sont $2-2\cos(2\pi k/n)$, doubles pour
$0<k<n/2$, d'espaces propres engendrés par $x\mapsto\cos(2\pi kx/n)$ et $x\mapsto\sin(2\pi kx/n)$.
La rotation $\varphi(x)=x+1$ est un automorphisme. Avec $\theta=2\pi k/n$,
\[ \cos(\theta(x+1))=\cos\theta\cos\theta x-\sin\theta\sin\theta x,\qquad
   \sin(\theta(x+1))=\sin\theta\cos\theta x+\cos\theta\sin\theta x, \]
donc, dans ces bases, $C$ est diagonale \emph{par blocs}, chaque bloc étant la rotation d'angle
$\theta$. Une correspondance qui conserve exactement la structure n'a donc pas une matrice
diagonale ; elle a une matrice qui respecte les espaces propres.
\end{exemple}

\begin{remarque}[Instabilité des modes proches]
\label{rem:daviskahan}
Si deux valeurs propres sont égales, le choix des vecteurs dans leur plan est arbitraire ; si elles
sont proches, un petit bruit les fait tourner dans ce plan. Le théorème de Davis et Kahan
\cite{daviskahan1970} borne l'angle entre un \emph{sous-espace} propre et son perturbé par la
perturbation divisée par l'écart qui sépare ce sous-espace du reste du spectre : les sous-espaces
bien séparés sont stables, les vecteurs individuels d'un groupe de valeurs proches ne le sont pas.
Toute comparaison mode à mode entre deux domaines doit donc regrouper les modes proches.
\end{remarque}

\subsection{La \enquote{dimension cachée} : l'analogie des corrélations canoniques}

L'analyse des corrélations canoniques \cite{hotelling1936} cherche, pour deux jeux de variables
mesurés sur les mêmes individus, les combinaisons $a^\top X$ et $b^\top Y$ de corrélation maximale.
Elle peut faire ressortir une composante \emph{faible dans chaque jeu pris seul}, que l'analyse en
composantes principales de chaque jeu ne voit pas.

\begin{exemple}[Une composante faible partout, parfaitement partagée]
\label{ex:cca}
Soient $G,H,Z$ indépendantes, centrées, de variance 1, et $0<\epsilon<1$. Posons $X=(G,\epsilon Z)$ et
$Y=(H,\epsilon Z)$. La première composante principale de $X$ est $G$ (variance $1>\epsilon^2$), non
corrélée à $Y$. Or
\[ \operatorname{corr}(a^\top X,b^\top Y)=\frac{a_2b_2\epsilon^2}{\sqrt{(a_1^2+a_2^2\epsilon^2)(b_1^2+b_2^2\epsilon^2)}}\le1, \]
avec égalité si et seulement si $a_1=b_1=0$ : la première paire canonique est $(Z,Z)$, de
corrélation 1.
\end{exemple}
\begin{proof}[Vérification]
La covariance croisée vaut $a_2b_2\epsilon^2$, et
\[ |a_2b_2|\epsilon^2=\sqrt{a_2^2\epsilon^2}\,\sqrt{b_2^2\epsilon^2}\le\sqrt{(a_1^2+a_2^2\epsilon^2)(b_1^2+b_2^2\epsilon^2)}, \]
l'inégalité étant stricte dès que $a_1\neq0$ ou $b_1\neq0$ (si $a_2b_2=0$, la corrélation est nulle).
\end{proof}

\textbf{L'analogie, dite comme telle} \cite[correction (4)]{complement89}. Les \enquote{individus} sont
les nœuds appariés par $\varphi$ ; les \enquote{variables} sont les coordonnées spectrales de chaque
graphe. La \enquote{dimension cachée} serait une paire de combinaisons de modes, peu marquée dans
chaque graphe seul (pas parmi ses modes dominants), mais alignée par $\varphi$ :
$u_A\approx Pu_B$. Seule $\varphi$ la révèle, comme seul l'appariement des individus révèle $Z$ dans
l'exemple~\ref{ex:cca}. Ce n'est qu'une analogie : aucun résultat de l'analyse des corrélations
canoniques ne se transporte ici sans démonstration.

\subsection{Où l'analogie casse}
\begin{enumerate}
\item \textbf{Graphes typés et orientés.} Le laplacien ci-dessus exige une matrice symétrique et
positive. Un graphe orienté a un laplacien de Chung, défini pour un graphe fortement connexe à
partir du vecteur de Perron d'une marche aléatoire \cite{chung2005} ; un graphe signé (activation,
inhibition) a un laplacien signé \cite{kunegis2010} ; pour des relations de plusieurs types, il
faut choisir une combinaison des types ou un opérateur par type. Aucun de ces choix n'est canonique,
et chacun change les modes.
\item \textbf{Modes proches} (exemple~\ref{ex:cycle}, remarque~\ref{rem:daviskahan}) : comparer des
vecteurs propres un à un est instable ; comparer des sous-espaces séparés l'est moins.
\item \textbf{Domaines éloignés.} Deux domaines éloignés ne sont pas presque isométriques : c'est la
définition même du mariage éloigné \cite[section~7]{q5}. Le défaut de commutation $\varepsilon$ de la
proposition~\ref{prop:fmapproche} n'a aucune raison d'être petit.
\item \textbf{Les premiers modes codent la surface.} Les modes grossiers d'un graphe de concepts
codent communautés et popularité, comme les premières composantes des plongements de mots codent
la fréquence \cite[p.~4]{mu2018} : un alignement du grossier au fin aligne d'abord la surface
\cite[section~7]{q5}. Il faut des opérateurs typés, normalisés par le degré, et le nul de $Z$.
\item \textbf{Tailles différentes, correspondances partielles.} $P$ n'est plus une permutation :
$C$ perd l'orthogonalité, la proposition~\ref{prop:fmapiso} ne s'applique plus telle quelle.
\end{enumerate}

\begin{statuts}[Section 4 : statuts]
\demontre Sous isomorphisme, $C$ commute avec les spectres et est diagonale si le spectre est simple
(\ref{prop:fmapiso}) ; sous presque isomorphisme, la borne $\varepsilon/|\lambda^A_i-\lambda^B_j|$
(\ref{prop:fmapproche}) ; le cycle donne une matrice par blocs (\ref{ex:cycle}) ; une composante
faible peut être parfaitement partagée (\ref{ex:cca}). Résultats élémentaires ; la remarque sur
Davis et Kahan est citée, non redémontrée.
\mesure Rien sur des graphes de domaines.
\conjecture Qu'une \enquote{dimension cachée} commune existe en général entre deux domaines qui ont
donné une découverte, et qu'un appariement spectral grossier-fin, typé et normalisé, la retrouve
plus souvent qu'un appariement paire à paire, à coût égal, sur un banc de correspondances plantées
(section~\ref{sec:recherche}).
\end{statuts}

%======================================================================
\section{Surprise, intérêt, portée et profondeur}
\label{sec:surprise}

\begin{definition}[Surprise de surface]
$S_t(a,b)=-\log p^0_t(a,b)$, où $p^0_t$ est le modèle de surface appris avant $t$
\cite[partie~2]{fable}.
\end{definition}

\begin{definition}[Crédibilité et intérêt]
\label{def:interet}
La crédibilité est $q=\sigma(q_0+\beta'Z-\vartheta S)$, $\sigma$ la fonction logistique,
$\beta',\vartheta>0$ ; l'intérêt est $I=q\,(S+W)$, gain d'information espéré, où $W$ est la portée
(définition~\ref{def:portee}) \cite[partie~2]{fable}. Absurde quand $q\to0$, évident quand $S\to0$.
\end{definition}

\begin{proposition}[La distance optimale recule quand la structure s'emboîte]
\label{prop:optimum}
Supposons $W$ indépendant de $S$, et $\partial W/\partial Z<\beta'q(S+W)$ (en particulier $W$
indépendant de $Z$). Posons $g(S,Z)=\vartheta(1-q)(S+W)-1$. Alors :
(i) $\partial I/\partial S=-q\,g$ ;
(ii) $g$ est strictement croissante en $S$, donc $I$ a au plus un maximum intérieur $S^\ast$, défini
par $\vartheta(S^\ast+W)(1-q)=1$, et c'est un maximum ;
(iii) là où il existe, $dS^\ast/dZ>0$.
\end{proposition}
\begin{proof}
(i) $\partial q/\partial S=-\vartheta q(1-q)$, donc $\partial I/\partial S=q-\vartheta q(1-q)(S+W)=-q\,g$.
(ii) $\partial g/\partial S=\vartheta\bigl[(1-q)+\vartheta q(1-q)(S+W)\bigr]>0$. Au zéro de $g$,
$\partial^2I/\partial S^2=-q\,\partial g/\partial S<0$. Si $g(0,Z)>0$ (portée très grande), $I$
décroît partout et il n'y a pas d'optimum intérieur.
(iii) $\partial q/\partial Z=\beta'q(1-q)$, donc
$\partial g/\partial Z=\vartheta(1-q)\bigl[\partial W/\partial Z-\beta'q(S+W)\bigr]<0$ par hypothèse ;
le théorème des fonctions implicites donne $dS^\ast/dZ=-g_Z/g_S>0$.
\end{proof}
La question~5 a refait cette dérivation et relevé que l'hypothèse \enquote{$W$ indépendant de $S$}
n'est pas dite chez Fable \cite[section~4, point~5]{q5} ; la condition sur $\partial W/\partial Z$
est ajoutée ici. Le score de Carbon, $I=d\cdot K$, n'a pas d'optimum intérieur. Épreuve : notation
à l'aveugle, plan $2\times2$ (surprise forte ou faible, serrure forte ou faible) \cite[section~4]{q7}.

\begin{definition}[Portée : surprise retirée ailleurs]
\label{def:portee}
Pour un lecteur de loi $P$ (modèle du public, lecteur figé, ou loi d'un monde planté), une idée $x$
et des faits $Y=(y_1,\dots,y_n)$ \emph{hors de son voisinage},
\[ W(x;Y)=\log P(Y\mid x)-\log P(Y), \]
information mutuelle ponctuelle \cite[section~2]{q2}. \enquote{Ailleurs} exclut le lien et ses
parents, les faits qui nomment un parent ou un voisin immédiat, ceux qui partagent avec le texte de
l'idée des entités au-delà d'un seuil fixé d'avance, et ceux qui ont servi à produire, choisir ou
valider l'idée.
\end{definition}

\begin{proposition}[Trois propriétés de la portée]
\label{prop:portee}
Sous une même loi $P$, avec $S(x)=-\log P(x)$ et $S(x\mid Y)=-\log P(x\mid Y)$ :
\begin{enumerate}[label=(\roman*)]
\item \textbf{symétrie} : $W(x;Y)=S(x)-S(x\mid Y)$ ;
\item \textbf{borne} : si $P$ est une loi de probabilité sur un ensemble discret, $W(x;Y)\le S(x)$ ;
\item \textbf{décomposition} : $W(x;Y)=\sum_jW(x;y_j)+U(x;Y)$, où
\[ U(x;Y)=\Bigl[\log P(Y\mid x)-\textstyle\sum_j\log P(y_j\mid x)\Bigr]-\Bigl[\log P(Y)-\textstyle\sum_j\log P(y_j)\Bigr] \]
est ce que l'idée ajoute à la dépendance entre les faits (l'unification de Myrvold \cite{myrvold2003}).
\end{enumerate}
\end{proposition}
\begin{proof}
(i) Par la formule de Bayes, $\log P(Y\mid x)-\log P(Y)=\log P(x\mid Y)-\log P(x)$.
(ii) $P(x\mid Y)\le1$, donc $S(x\mid Y)\ge0$ et $W=S(x)-S(x\mid Y)\le S(x)$.
(iii) Ajouter et retrancher $\sum_j\log P(y_j\mid x)$ et $\sum_j\log P(y_j)$ ; le reste est
$\sum_jW(x;y_j)$.
\end{proof}

\begin{exemple}[Une unification sans aucune prédiction lien par lien]
\label{ex:unification}
Deux faits binaires $y_1,y_2$, indépendants et équiprobables sous $P$. Sous l'idée $x$, chacun reste
équiprobable mais $y_1=y_2$. On observe $Y=(1,1)$ : $W(x;y_j)=0$ pour chaque fait, mais
$P(Y)=1/4$, $P(Y\mid x)=1/2$, d'où $W(x;Y)=U=\log2>0$. Une somme lien par lien, comme la formule de
Fable \cite[partie~2]{fable}, ne voit pas ce gain \cite[section~2]{q2}.
\end{exemple}

\begin{remarque}[Bruit et paraphrase ; pondération]
D'après l'erratum de la question~2 \cite{erratamineurs} : \emph{$W$, borné par $S$, écarte la
paraphrase} (peu surprenante, elle ne retire presque rien) ; \emph{le rapport $W/S$ écarte le bruit}
(très surprenant, il ne retire rien). La borne (ii) vaut pour un lecteur cohérent (monde planté),
non pour un lecteur figé. Pondérer $W$ par la centralité des faits fait revenir le degré des parents
par construction : $W$ principal non pondéré, version pondérée en second, avec sa corrélation
partielle au degré \cite[section~2]{q2}.
\end{remarque}

\begin{test}[La portée explique-t-elle l'effet du nœud chez Gu et Krenn ?]
\label{test:gukrenn}
\emph{Fait de départ} : sur 2\,996 idées bâties sur une paire de concepts (4\,451 notes, 110
responsables d'équipe), l'intérêt décroît avec le degré du concept de l'évaluateur,
$\rho=-0{,}166$ (PageRank $-0{,}176$, distance $-0{,}223$) \cite[v4, p.~8]{gukrenn2026} ; c'est un
fait sur le \emph{parent}, non sur l'idée fille, et il ne pèse qu'environ 3~\% de la variance des
rangs \cite[section~1]{q2}. \emph{Statistique} : corrélation partielle de rang entre $W$ estimé
(lecteur figé, cent titres de 2024-2025 hors voisinage, faux appariés) et la note, à degré,
PageRank, citations, distance, longueur et évaluateur fixés. \emph{Nul} : réaffectations
idée-faits permutées. \emph{Prédiction} : au moins $+0{,}05$, borne basse à 95~\% (rééchantillonnage
par évaluateur) au-dessus de 0, et le coefficient du degré perd au moins un quart. \emph{Réfutation} :
borne haute sous $+0{,}02$, gardes tenues ; indécidable si $W$ estimé est corrélé au degré au-delà
de 0,5 \cite[section~6]{q2}. Demi-largeur sans grappes environ 0,036 : le seuil frôle la limite,
puissance à calculer au préenregistrement.
\end{test}

\begin{statuts}[Section 5 : statuts]
\demontre L'optimum de surprise et son recul quand $Z$ croît, sous des hypothèses dites
(\ref{prop:optimum}) ; symétrie, borne et décomposition de la portée (\ref{prop:portee}) ;
l'exemple d'unification pure (\ref{ex:unification}).
\mesure Les chiffres de Gu et Krenn, lus à la page par la question~2 et confirmés par la
contre-lecture \cite[section~Q2]{contrelecture}.
\conjecture Que $W$ estimé explique l'intérêt au-delà du nœud (épreuve~\ref{test:gukrenn}) ; que
l'optimum de distance soit mobile (Fable) plutôt que fixe (Carbon) \cite[section~4]{q7}.
Dans l'organisme d'idéation, $W$ estimé n'est qu'une mesure rapportée tant qu'un monde planté à
portée connue ne l'a pas validé ; il n'entre jamais dans un choix de proposeur ni dans une
récompense (règle de la décision N-053, point~5, citée par la question~2).
\end{statuts}

%======================================================================
\section{Le puits profond et étroit, et la fécondité}
\label{sec:puits}

Le complément du 24 novembre demande trois grandeurs, distinctes, avec pour chacune une définition
formelle et une épreuve falsifiable sur données datées, avec son nul
\cite{complementspectral}. Pour l'organisme : la profondeur comme récompense, l'étroitesse comme
difficulté, la fécondité comme critère différé.

\subsection{Profondeur d'un puits}

Le complément définit la profondeur comme la surprise retirée ailleurs, réalisée par trois
canaux : les relations transportées, les faits importés comme prédictions, la stabilité sous de
petites perturbations \cite[point (2)]{complementspectral}. La question~2 réserve
\enquote{ailleurs} aux faits hors du voisinage. Les deux ne coïncident pas : les prédictions de
$\varphi$ (cas C3 et C4) tombent \emph{dans} le voisinage de $b$. La note garde donc deux échelles,
nommées et rapportées à part.

\begin{definition}[Profondeur locale et profondeur à distance]
\label{def:profondeur}
Pour une correspondance $\varphi$ gelée à la date $t$ et un horizon $h$ :
\begin{enumerate}
\item \textbf{transport} : $Z_r(\varphi)$, le score $c-\kappa$ réduit sous recâblage ;
\item \textbf{importations} : l'ensemble $\Pi_t(\varphi)$ des prédictions de $\varphi$ (relations
absentes, cas C3 ; profils des objets manquants, cas C4), et leur \emph{taux de confirmation}
$\pi_{t,h}(\varphi)$, part de $\Pi_t(\varphi)$ affirmée dans $(t,t+h]$ ;
\item \textbf{robustesse} : $\rho_\varepsilon(\varphi)$, probabilité que $c-\kappa$ reste à $\tau$
près sous une perturbation des faits de taux $\varepsilon$ ;
\item \textbf{portée} : $W$ sur des faits hors du voisinage (définition~\ref{def:portee}).
\end{enumerate}
La \emph{profondeur locale} est le triplet $(Z_r,\pi_{t,h},\rho_\varepsilon)$ ; la \emph{profondeur à
distance} est $W$. Aucune combinaison n'est fixée ici ; elle se fixe avant un préenregistrement.
\end{definition}

\begin{test}[Les prédictions d'une correspondance profonde se confirment-elles ?]
\label{test:profondeur}
\emph{Données} : relations typées et datées de PubTator~3 \cite{pubtator2024}, gel fin 2010
(l'expérience de Fable \cite[partie~7]{fable}). \emph{Unités} : paires $(a,b)$ avec leur meilleure
correspondance $\varphi$ (criblage par signatures de rôle, puis transport sur liste courte) et ses
prédictions $\Pi_{2010}(\varphi)$. \emph{Issue} : première affirmation d'une relation prédite en
2011-2015. \emph{Statistique} : rapport des taux de confirmation des prédictions des correspondances
du tiers supérieur de profondeur locale, contre deux témoins appariés sur la popularité des
extrémités et les types : (i) prédictions de correspondances tirées sous recâblage (même degrés,
mêmes types) ; (ii) prédictions des correspondances voisines (une transposition) des mêmes paires.
\emph{Prédiction} (seuil proposé, à geler) : rapport au moins 1,5, borne basse à 95~\% au-dessus de 1
(rééchantillonnage par concept). \emph{Réfutation} : borne haute sous 1,1, gardes tenues.
\emph{Gardes} : sur un monde planté (correspondance et prédictions plantées, dévoilées après le gel),
le rapport doit être grand ; après permutation des types à degrés constants, il doit tomber à 1,
sinon la mesure lit la popularité. \emph{Volet robustesse} : à transport égal, les prédictions des
correspondances robustes se confirment plus que celles des fragiles.
\end{test}

\subsection{Étroitesse}

\begin{definition}[Étroitesse]
\label{def:etroitesse}
L'\emph{étroitesse} de $\varphi$ est $\mathrm{Etr}_\tau(\varphi)=1-L_\tau(\varphi;1)$, complétée par
$1-L_\tau(\varphi;2)$ (définition~\ref{def:largeur}) : la petite part des correspondances voisines
qui marchent encore \cite[point (2)]{complementspectral}. On la rapporte avec la rareté au hasard
$\mathrm{Rar}_\tau(\varphi)$, qui est une autre grandeur (section~\ref{sec:perturbations}).
\end{definition}

Le couple \enquote{profond et étroit} est l'aiguille : un puits que peu de tirages touchent et
qu'une petite erreur de correspondance fait perdre. D'où l'intérêt de partir des anomalies, qui
réduisent l'espace à explorer (section~\ref{sec:recherche}).

\begin{test}[L'étroitesse retarde-t-elle la découverte ?]
\label{test:etroitesse}
\emph{Données} : les naissances de 2011-2020 dont les deux voisinages typés ont atteint une taille
fixée d'avance à une date de maturité $t_0$ antérieure. \emph{Issue} : délai entre $t_0$ et la
naissance. \emph{Statistique} : modèle de Cox avec étroitesse, profondeur locale, $S$, $F$, $P$,
erreurs groupées par concept. \emph{Nul} : permutation de l'étroitesse dans des strates de
profondeur et de surprise. \emph{Prédiction} : rapport de risques par écart type d'étroitesse
inférieur à 1 (les puits étroits naissent plus tard), seuil proposé 0,9, borne haute sous 1.
\emph{Réfutation} : borne basse au-dessus de 0,95, gardes tenues. \emph{Version dans l'organisme} :
sur un banc de correspondances plantées où l'étroitesse est plantée (leurres proches de $\varphi$) à
rareté fixée, le coût en jetons pour retrouver $\varphi$ doit croître avec l'étroitesse ; à
étroitesse fixée, avec la rareté.
\end{test}

\subsection{Fécondité}
\label{sec:fecondite}

La fécondité est ce que l'idée rend possible après elle : en cinétique, un catalyseur qui abaisse
la barrière d'autres réactions \cite[point (3)]{complementspectral} (section~\ref{sec:cinetique}).

\begin{definition}[Fécondité]
\label{def:fecondite}
Pour une naissance $(a,b)$ à la date $t_b$, de correspondance $\varphi$, et un horizon $h$ :
\begin{enumerate}
\item \textbf{gain de prédiction des liens futurs} : $\Phi_h=\mathrm{PM}_1-\mathrm{PM}_0$, où
$\mathrm{PM}$ est la précision moyenne, sur l'ensemble à risque, de la prédiction des liens nés dans
$(t_b,t_b+h]$, par un modèle de base ($\mathrm{PM}_0$) et par le même modèle augmenté du trait
\enquote{paire alignée par $\varphi$} ($\mathrm{PM}_1$) : c'est le \enquote{cortège} de Fable
\cite[partie~7, C]{fable} ;
\item \textbf{indice de rupture} de l'article de naissance : avec $N_A$ les travaux ultérieurs qui
citent l'article sans citer ses références, $N_B$ ceux qui citent l'article et ses références,
$N_C$ ceux qui ne citent que ses références,
$\mathcal{D}=(N_A-N_B)/(N_A+N_B+N_C)\in[-1,1]$ \cite[p.~7, légende de la figure~1 ; p.~18]{wu2019}.
C'est la forme employée par Wu, Wang et Evans ; l'indice de Funk et Owen-Smith \cite{funk2017}, dont
il dérive, est lu sur notice et résumé, sa définition exacte est à vérifier sur la source.
\end{enumerate}
\end{definition}

\begin{remarque}[Profondeur et fécondité sont distinctes]
Une unification qui ne prédit rien de neuf (toutes ses conséquences déjà connues) a une portée
$W$ réalisée nulle sur l'avenir mais peut avoir un fort transport ; un instrument nouveau, sans
aucune correspondance, peut multiplier les naissances voisines (fécondité forte, profondeur nulle).
Les deux grandeurs se mesurent donc séparément. La validité de l'indice de rupture contre le
jugement des pairs reste discutée \cite{bornmann2020,leibel2024} ; c'est une covariable, non un
critère \cite[section~2]{q2}. Une baisse séculaire de la rupture est rapportée par Park, Leahey et
Funk \cite{park2023} (titre et notice) : toute comparaison entre périodes doit en tenir compte.
\end{remarque}

\begin{test}[La profondeur mesurée à la naissance prédit-elle la fécondité ?]
\label{test:fecondite}
\emph{Données} : naissances de 2011-2015 de l'expérience de Fable, liens de 2016-2025.
\emph{Statistiques} : (i) cortège : rapport des taux de naissance des paires alignées par $\varphi$
contre des paires croisées non alignées, appariées (types, popularité, proximité de surface) ;
prédiction de Fable au moins 1,5, borne basse au-dessus de 1 \cite[partie~7]{fable} ; (ii) régression
de $\Phi_{10}$ et de $\mathcal{D}$ sur la profondeur locale à $t_b$, à $S$, $Z$, $F$, $P$ fixés.
\emph{Nul} : placebo temporel, la même mesure sur la fenêtre \emph{précédant} la naissance doit
donner zéro ; permutation des naissances entre paires de même strate. \emph{Réfutation} : cortège à
borne haute sous 1,1 ; coefficient de la profondeur à borne haute sous 0, gardes tenues. Les deux
issues instruisent : si la profondeur ne prédit pas la fécondité, les deux sont confirmées
distinctes.
\end{test}

\subsection{Usage dans l'organisme d'idéation}
\begin{itemize}
\item \textbf{Profondeur comme récompense} : seulement là où elle est connue par construction, dans
un monde planté dont l'orchestrateur seul tient la graine ; la composition des proposeurs se juge
sur l'oracle du monde planté et le coût, jamais sur un juge (règle de la décision N-058, rappelée
par la question~2) ; la profondeur \emph{estimée} reste une mesure rapportée.
\item \textbf{Étroitesse comme difficulté} : elle fixe le budget d'exploration et le choix des
méthodes de recherche (section~\ref{sec:recherche}), non la note d'une fiche.
\item \textbf{Fécondité comme critère différé} : jugée sur des liens apparus après coup ; dans un
monde planté, sur des faits cachés dévoilés plus tard par l'orchestrateur.
\item \textbf{Gonflement} : un proposeur peut gonfler le transport en ajoutant des objets
(proposition~\ref{prop:combler}), l'étroitesse en choisissant des leurres, la portée en amorçant la
batterie de faits \cite[section~5]{q2}. Parades : coût de description des objets ajoutés, batterie
scellée et renouvelée, valeurs jamais rendues aux proposeurs.
\end{itemize}

\begin{statuts}[Section 6 : statuts]
\demontre Rien de neuf : la section définit et renvoie aux propositions des sections 3 et 5.
\mesure La forme de l'indice de rupture, lue à la page \cite{wu2019}.
\conjecture Les épreuves~\ref{test:profondeur}, \ref{test:etroitesse} et \ref{test:fecondite}, avec
leurs seuils proposés, à geler.
\end{statuts}

%======================================================================
\section{Le test qui reste sur la résonance}
\label{sec:resonance}

Fable écarte la résonance comme mécanisme et garde un test : la co-accélération d'activité des deux
voisinages prédit-elle la naissance au-delà du champ $F$ ? \cite[partie~2]{fable}. Carbon l'écarte
explicitement, faute de phase mesurable \cite[l.~64]{carbon}. La note écrit le test.

\begin{definition}[Activité, croissance, co-accélération]
\label{def:coacc}
Pour un concept $x$, soit $n_x(t)$ le nombre de documents de l'année $t$ qui emploient un concept du
voisinage $N(x)$ (rayon 1, dans la couche d'usage). Croissance : $g_x(t)=\Delta\log(1+n_x(t))$ ;
accélération : $\alpha_x(t)=g_x(t)-g_x(t-1)$, centrée réduite par champ et année. La
\emph{co-accélération} de $(a,b)$ sur une fenêtre de $w$ ans finissant à $t-\ell$ (décalage $\ell\ge1$
pour exclure l'anticipation de la découverte) est
\[ \mathrm{CoA}_{ab}(t)=\frac1w\sum_{s=t-\ell-w+1}^{t-\ell}\alpha_a(s)\,\alpha_b(s). \]
\end{definition}

\begin{test}[Résonance]
\label{test:resonance}
\emph{Énoncé} : dans le modèle de Cox
\[ \log\lambda_{ab}(t)=\log\lambda_0(t)+\log p^0_t(a,b)+\beta Z_{ab}+\gamma F_t(a,b)+\delta P_t(a,b)+\zeta\,\mathrm{CoA}_{ab}(t), \]
$\zeta>0$. \emph{Données} : l'ensemble à risque de l'expérience de Fable, naissances 2011-2015, gel
fin 2010 pour $Z$, séries d'activité annuelles avant chaque date. \emph{Statistique} : rapport de
risques par écart type de $\mathrm{CoA}$, intervalle à 95~\% par rééchantillonnage par concept ;
rapport de vraisemblance pour $\zeta$. \emph{Nuls} : (i) décalage circulaire, tiré au hasard, de la
série de $b$ (conserve les accélérations marginales, brise la synchronie) ; (ii) remplacement de
$b$ par un concept $b'$ de même champ et de même accélération moyenne. \emph{Contrôles} : $F$ (les
cadres communs accélèrent ensemble les deux voisinages), $P$ (une cible commune), effets fixes
champ-année (un choc commun), accélérations marginales $\alpha_a$, $\alpha_b$ seules.
\emph{Prédiction} (seuil proposé, à geler) : rapport de risques au moins 1,10, borne basse au-dessus
de 1, et effet réel supérieur à 95~\% des effets sous le nul (i). \emph{Réfutation} : borne haute sous
1,05, gardes tenues ; ou effet sous nul aussi fort que l'effet réel ; ou effet absorbé par $F$.
\emph{Garde} : un monde planté où deux voisinages co-accélèrent sans pont (le témoin doit mourir) et
un autre où la co-accélération est plantée comme cause.
\end{test}

\begin{statuts}[Section 7 : statuts]
\demontre Rien.
\mesure Rien.
\conjecture $\zeta>0$ au-delà de $F$ ; à défaut, la \enquote{résonance} n'a aucun contenu propre et
le mot doit sortir du modèle.
\end{statuts}

%======================================================================
\section{Cinétique et émergence}
\label{sec:cinetique}

\subsection{La cinétique des collisions d'idées}

La question~3 relit le taux fusionné comme une cinétique chimique \cite{q3}, et son erratum en
corrige les budgets \cite{erratumq3}. Le taux~\eqref{eq:taux} est déjà une loi d'Arrhenius
$k=A\,e^{-E/\Theta}$ \cite{arrhenius1889}, avec $E/\Theta=S-\beta Z-\gamma F-\delta P$ et l'unité
d'énergie fixée par le coefficient 1 de $S$ (un nat de surprise du public).

\begin{table}[ht]
\centering\small
\caption{Dictionnaire de la cinétique, abrégé d'après la question~3, section~1 \cite{q3}.}
\label{tab:dico}
\begin{tabularx}{\textwidth}{@{}lXl@{}}
\toprule
Terme & Grandeur mesurable & Statut \\
\midrule
Concentration & $n_a(t)$ : auteurs actifs sur $a$ ; dans l'organisme, contextes où $a$ figure & ajoute (ordre) \\
Collision & rencontre sans prise (coprésence sans relation extraite) & ajoute \\
Facteur d'orientation & $e^{\beta Z}$ & redonne $Z$ \\
Énergie d'activation & corpus : seulement $E/\Theta$ ; organisme : pente d'Arrhenius & organisme seulement \\
Température & organisme : température d'échantillonnage (forme de Boltzmann \cite[éq.~(1), p.~2-3]{hinton2015}) & ajoute dans l'organisme \\
Catalyseur & méthode présente, non consommée, ouvrant une voie parallèle & ajoute ; $F$ en est un \\
Champ physique & déplacement uniforme des barrières & écarté (taux de base) \\
Autocatalyse & multiplicateur des paires alignées, rapport de branchement & redonne ; précurseur \\
\bottomrule
\end{tabularx}
\end{table}

\begin{figure}[ht]
\centering
\begin{tikzpicture}[x=1cm,y=1cm,st/.style={draw,rounded corners,inner sep=4pt,font=\small}]
\node[st] (s0) at (0,0) {séparés $a+b$};
\node[st] (s1) at (4.2,0) {rencontre $\{a,b\}$};
\node[st] (s2) at (8.4,0) {idée née};
\draw[-{Stealth[length=5pt]}] ($(s0.east)+(0,0.12)$) -- node[above,font=\small]{$k_D$} ($(s1.west)+(0,0.12)$);
\draw[-{Stealth[length=5pt]}] ($(s1.west)+(0,-0.12)$) -- node[below,font=\small]{$k_{\mathrm{sep}}$} ($(s0.east)+(0,-0.12)$);
\draw[-{Stealth[length=5pt]}] (s1) -- node[above,font=\small]{$k_r$} (s2);
\end{tikzpicture}
\caption{Le schéma \enquote{rencontre, séparation, prise} de la question~3, prédiction~1.}
\label{fig:schema}
\end{figure}

\begin{proposition}[Temps moyen et étape limitante]
\label{prop:tempsmoyen}
Dans la chaîne de Markov de la figure~\ref{fig:schema} (taux constants), le temps moyen jusqu'à la
naissance depuis l'état séparé est
\[ \frac1\lambda=\frac1{k_D}+\frac1{k_r}+\frac{k_{\mathrm{sep}}}{k_D\,k_r}. \]
Rencontres brèves ($k_{\mathrm{sep}}\gg k_D,k_r$) : $\lambda\approx(k_D/k_{\mathrm{sep}})\,k_r$, un
produit, c'est le taux fusionné. Rencontres durables ($k_{\mathrm{sep}}\to0$) :
$\lambda=k_Dk_r/(k_D+k_r)$, forme en série \cite{collins1949} (page à vérifier).
\end{proposition}
\begin{proof}
Soient $T_0$ et $T_1$ les temps moyens depuis les états séparé et rencontre. Par l'analyse au
premier pas,
\[ T_0=\frac1{k_D}+T_1,\qquad T_1=\frac1{k_{\mathrm{sep}}+k_r}+\frac{k_{\mathrm{sep}}}{k_{\mathrm{sep}}+k_r}\,T_0 . \]
D'où
\[ T_0\,\frac{k_r}{k_{\mathrm{sep}}+k_r}=\frac1{k_D}+\frac1{k_{\mathrm{sep}}+k_r},\qquad
   T_0=\frac{k_{\mathrm{sep}}+k_r}{k_Dk_r}+\frac1{k_r}, \]
qui est la formule annoncée. Les deux limites s'en lisent.
\end{proof}

\begin{proposition}[Le signe de l'interaction dit le régime]
\label{prop:interaction}
Dans la forme générale,
\[ \frac{\partial^2\log\lambda}{\partial\log k_D\,\partial\log k_r}=\frac{k_Dk_r}{(k_D+k_r+k_{\mathrm{sep}})^2}>0 . \]
En régime durable, l'élasticité au contact vaut $k_r/(k_D+k_r)$, celle de la prise $k_D/(k_D+k_r)$,
de somme 1.
\end{proposition}
\begin{proof}
$\lambda=k_Dk_r/(k_D+k_r+k_{\mathrm{sep}})$, donc $\log\lambda=\log k_D+\log k_r-\log\Sigma$ avec
$\Sigma=k_D+k_r+k_{\mathrm{sep}}$. Alors $\partial\log\lambda/\partial\log k_D=(k_r+k_{\mathrm{sep}})/\Sigma$,
puis $\partial/\partial\log k_r$ de cette expression vaut
$[k_r\Sigma-(k_r+k_{\mathrm{sep}})k_r]/\Sigma^2=k_Dk_r/\Sigma^2$. Avec $k_{\mathrm{sep}}=0$ on lit les
élasticités.
\end{proof}
Un signe négatif de l'interaction contact $\times$ structure réfute donc toute la famille ; un signe
nul dit que la cinétique n'ajoute rien au taux fusionné sur ce point \cite[section~2]{q3}. Dans la
simulation du régime durable (partie~B, non touchée par l'erratum), l'interaction vaut $+0{,}118$
(écart type 0,006), l'élasticité au contact $0{,}30$, $0{,}42$, $0{,}55$ du tiers bas au tiers haut de
$Z$ \cite[section~2]{q3} (sortie \texttt{500e94df…}).

\begin{proposition}[Le terme des cadres est le premier ordre d'un catalyseur]
\label{prop:catalyseur}
Soit la voie catalytique de la question~3 corrigée par son erratum \cite[mineur~2]{erratumq3} :
\[ \lambda=p^0\Bigl[e^{\beta Z+\delta P}+\psi\sum_mw_m\,h_a(m)h_b(m)\,e^{\beta_cZ+\delta P}\Bigr]. \]
Avec $\tilde F=\sum_mw_mh_a(m)h_b(m)$ :
$\log\lambda=\log p^0+\beta Z+\delta P+\log\bigl(1+\psi\tilde F e^{(\beta_c-\beta)Z}\bigr)$. Au premier
ordre en $\psi\tilde F$, on retrouve $\gamma F$ avec $\gamma=\psi$, $F=\tilde F$, si et seulement si
$\beta_c=\beta$. De plus $\log\lambda$ est concave en $\tilde F$, et si $\beta_c<\beta$ le gain relatif
$\psi\tilde Fe^{(\beta_c-\beta)Z}$ décroît avec $Z$.
\end{proposition}
\begin{proof}
Factoriser $e^{\beta Z+\delta P}$ ; $\log(1+x)=x+O(x^2)$ ; le terme d'ordre un vaut
$\psi\tilde Fe^{(\beta_c-\beta)Z}$, indépendant de $Z$ si et seulement si $\beta_c=\beta$ ;
$\log(1+x)$ est concave ; la dérivée en $Z$ du gain relatif est $(\beta_c-\beta)$ fois ce gain.
Sans le facteur $e^{\delta P}$ sur la voie catalytique, le terme d'ordre un garderait $e^{-\delta P}$.
\end{proof}

\begin{proposition}[À température constante, la barrière n'est pas identifiable]
\label{prop:identif}
Si $\lambda=\nu\,e^{-E/\Theta}$ avec $E=E_0-\beta_EZ-\delta_EP$, la transformation
$(E_0,\beta_E,\delta_E,\Theta)\mapsto(sE_0,s\beta_E,s\delta_E,s\Theta)$, $s>0$, laisse $\lambda$
inchangé : seuls les rapports $E_0/\Theta$, $\beta_E/\Theta$, $\delta_E/\Theta$ sont identifiables, et
$E_0/\Theta$ se confond avec $\log\nu$.
\end{proposition}
\begin{proof}
$E/\Theta$ est invariant par la transformation ; $\log\lambda=\log\nu-E_0/\Theta+\dots$ ne sépare pas
$\log\nu$ de $E_0/\Theta$.
\end{proof}
La barrière s'identifie, à une unité près, si une variable de strate multiplie tous les
coefficients de barrière par un même facteur $1/\Theta_s$ (matrice des coefficients par strate de
rang 1), ou dans l'organisme, où $\Theta=T$ est la température d'échantillonnage, fixée par
l'expérimentateur \cite[section~3]{q3}.

\begin{table}[ht]
\centering\small
\caption{Budgets corrigés de la mesure de la barrière dans l'organisme (barrière de 2,5 nats),
d'après l'erratum 1 de la question~3 et sa sortie (\texttt{dc98b08f…}) \cite{erratumq3}. Forme
posée : $\ln k=\ln A+m\ln T-E/T$, empruntée à Eyring, à éprouver.}
\label{tab:budgets}
\begin{tabular}{@{}lrrr@{}}
\toprule
 & 1\,600 collisions & 6\,400 & 25\,600 \\
\midrule
Ancien plan (4 températures, 0,4 à 1,3), sans le terme en $\ln T$ & 0,56 & -- & -- \\
Ancien plan, avec le terme & 2,96 & 1,48 & 0,74 \\
Plan large (6 températures, 0,25 à 2,0), avec le terme, $m=0$ & 1,87 & 0,93 & 0,47 \\
Plan large, avec le terme, $m=1$ & 2,25 & 1,12 & 0,56 \\
\bottomrule
\end{tabular}
\par\smallskip
\begin{minipage}{0.92\textwidth}\footnotesize
Écarts types asymptotiques de l'estimation de $E$. Corrélation entre $\ln T$ et $-1/T$ : 0,987
(ancien plan), 0,962 (plan large). Omettre le terme s'il existe ($m=1$) biaise la lecture d'environ un
nat (3,44 à 3,85 au lieu de 2,5), sans baisse avec le budget. Test du catalyseur contre le
préfacteur à 12\,800 collisions : puissance 0,81 à 0,99 selon le plan et $m$, faux positifs 0,02 à
0,06, si le catalyseur ne change pas la dépendance du préfacteur en température. Au moins cinq
températures, non trois.
\end{minipage}
\end{table}

\textbf{Les autres prédictions de forme} \cite[section~2]{q3}, toutes supposées :
\emph{ordre en porteurs} ($\lambda\propto n_a^{m_a}n_b^{m_b}n_{ab}^{m_{ab}}$, trois mécanismes, trois
signatures ; réfutée si les ordres sont non entiers et instables entre champs) ; \emph{précurseur
structural de l'avalanche} (la taille du plus grand ensemble réflexivement autocatalytique nourri
par les concepts établis \cite{hordijk2004} prédit-elle un rapport de branchement d'au moins 1 sur la
période suivante ? forme occupée, test prédictif à faire) ; \emph{réversibilité} (pente entre 0 et 1
entre taux de naissance et survie des ponts, selon Evans et Polanyi \cite{evanspolanyi1938}, page à
vérifier). La question~3 conclut \enquote{garder en partie}, confiance 0,6, verdict maintenu par
l'erratum, la barrière se mesurant dans l'organisme \enquote{avec environ seize à vingt-huit fois
plus de collisions que la v1 ne le disait} \cite{erratumq3}.

\subsection{L'émergence}

La question~4 sépare deux hypothèses de sens opposés sur le contact \cite[section~1]{q4} :
\textbf{(a)} la structure d'un champ (le taux de naissance change quand le graphe des coauteurs
acquiert une composante géante ; en 2009, la transition \emph{suit} l'invention selon ses auteurs
\cite[p.~5 et 14 du manuscrit]{bettencourt2009}) ; \textbf{(b)} le processus d'un petit groupe
(l'interaction pendant la génération réduit le nombre d'idées distinctes ; chez Diehl et Stroebe,
le groupe nominal est devant dans 18 expériences sur 22, la cause retenue est le blocage de
production \cite[p.~497]{diehlstroebe1987}, absent chez des agents qui écrivent en parallèle).
L'écriture $D=D_{\mathrm{seul}}+D_{\mathrm{composé}}$ les rend complémentaires sous condition :
$D_{\mathrm{seul}}$ décroît avec le contact si les producteurs convergent, $D_{\mathrm{composé}}$
(idées qui exigent des pièces détenues par d'autres) est nul sans contact et croît avec la
connectivité.

\begin{proposition}[Cascades sous-critiques]
\label{prop:branchement}
Si chaque naissance en induit un nombre aléatoire de moyenne $n<1$, indépendamment, la taille
moyenne totale d'une cascade issue d'une naissance est $1/(1-n)$ ; si $n>1$, la cascade est
infinie avec une probabilité positive.
\end{proposition}
\begin{proof}
La génération $k$ a une taille moyenne $n^k$ (récurrence et indépendance) ; la somme
$\sum_kn^k=1/(1-n)$. Le second point est le théorème classique d'extinction d'un processus de
Galton et Watson \cite[chapitre~I]{harris1963} (page à vérifier).
\end{proof}

\begin{proposition}[Dilution d'arêtes et seuil]
\label{prop:dilution}
Dans un graphe aléatoire où chaque arête possible est présente avec probabilité $c/(N-1)$
indépendamment, conserver chaque arête avec probabilité $f$ indépendamment donne un graphe de même
loi avec $cf/(N-1)$ : le degré moyen devient $cf$, et le seuil d'apparition d'une composante géante,
en taille infinie, est $cf=1$. Dans le graphe biparti de Fable (degrés croisés moyens $c_A$, $c_B$),
le processus d'exploration qui alterne les deux côtés a pour moyenne en deux pas $c_Ac_Bf^2$, d'où le
seuil $c_Ac_Bf^2=1$.
\end{proposition}
\begin{proof}
Une arête est présente après dilution si et seulement si elle l'était et a été conservée, avec
probabilité $\frac{c}{N-1}f$, indépendamment des autres. Le seuil $cf=1$ pour un graphe de degrés de
Poisson est classique \cite{molloyreed1995,newman2001pre} (résumés). Pour le biparti : depuis un
nœud de $A$, le nombre de voisins conservés dans $B$ a pour moyenne $c_Af$ ; depuis chacun, $c_Bf$ ;
le processus de branchement à deux étapes a pour moyenne $c_Ac_Bf^2$, et la proposition
\ref{prop:branchement} donne le seuil. Le passage du processus de branchement au graphe (graphe
localement arborescent) est admis, non redémontré ici ; dérivation de la question~4
\cite[section~4]{q4}.
\end{proof}

\textbf{Deux définitions du saut} \cite{erratumq4}. Dans l'épreuve de relais (chaînes de $L$
sauts), la lecture \enquote{arête} donne $r(L)=h\,f^L$ pour $L=0$ à 4 ($L=0$ : un agent tient les deux
moitiés), $h=r(0)$ ; la lecture \enquote{relais} donne $r(0)=h$ et $r(L)=hf^{L-1}$ pour $L\ge1$ (contact
direct sans perte). Le rapport $r(1)/r(0)$ les sépare : il vaut $f$ dans la première, 1 dans la
seconde ; à 50 essais par longueur, son écart type est 0,127 pour $f=0{,}8$, il faut environ 200 essais
à $L=0$ et $L=1$. La formule de la v1, $r(L)=hf^{L-1}$ appliquée avec la percolation diluée, aurait
pris $h=r(1)=hf$ et sous-estimé la courbe d'environ 20~\%.

\begin{table}[ht]
\centering\small
\caption{Courbes de dose prédites à 16 agents, $f=0{,}8$, d'après l'erratum 1 de la question~4 et
sa sortie (\texttt{a31ee52c…}) \cite{erratumq4}.}
\label{tab:dose}
\begin{tabular}{@{}lccc@{}}
\toprule
Degré moyen $c$ & 1 & 2 & 4 \\
\midrule
Lecture \enquote{arête}, percolation & 0,111 & 0,353 & 0,656 \\
Lecture \enquote{arête}, voisins directs seuls & 0,056 & 0,093 & 0,165 \\
Lecture \enquote{relais}, percolation & 0,129 & 0,382 & 0,660 \\
Lecture \enquote{relais}, voisins directs seuls & 0,065 & 0,111 & 0,201 \\
Formule fautive ($h=r(1)$) & 0,092 & 0,284 & 0,525 \\
\bottomrule
\end{tabular}
\end{table}

La prédiction de percolation est \enquote{sans paramètre ajusté sur la courbe de dose, sous quatre
hypothèses posées} : $f$ ne dépend pas du degré (d'où une épreuve de relais à la charge de chaque
dose), les sauts sont indépendants, la prise $h$ est la même après un relais, le graphe réel de
chaque course est connu \cite{erratumq4}. Si l'une tombe, c'est la courbe \enquote{percolation avec
$f$ tiré des chaînes} qui est réfutée, non l'hypothèse d'émergence. Trois courses par dose, à quatre
composites chacune, départagent percolation et voisins directs dans environ 98~\% des cas, dans les
deux lectures. Ordre recommandé : réplique de Bettencourt et collègues sur OpenAlex (processeur,
gratuite) ; épreuve de relais puis courbe de dose (environ 130 dollars aux tarifs posés, non
engagés) ; contraste séparés contre interaction (le plus cher) \cite[section~6]{q4}.

\begin{statuts}[Section 8 : statuts]
\demontre Temps moyen et limites (\ref{prop:tempsmoyen}) ; signe de l'interaction
(\ref{prop:interaction}) ; premier ordre du catalyseur (\ref{prop:catalyseur}) ; non-identifiabilité
à température constante (\ref{prop:identif}) ; taille moyenne des cascades (\ref{prop:branchement}) ;
dilution d'arêtes (\ref{prop:dilution}, le seuil étant cité).
\mesure Par simulation, sous paramètres posés : budgets (tableau~\ref{tab:budgets}), courbes de dose
(tableau~\ref{tab:dose}). Rien sur des données réelles.
\conjecture Que les naissances d'idées obéissent à une cinétique ; que la barrière plantée en nats
représente la difficulté d'une vraie idée ; que la percolation prédise le recouvrement des
composites dans l'organisme.
\end{statuts}

%======================================================================
\section{Chercher dans un puits étroit ; amender le Plan Turing}
\label{sec:recherche}

\subsection{Deux niveaux}
La question~5 sépare deux niveaux \cite[section~7]{q5}. \emph{Quelle paire, quelle idée essayer} :
réglé le 22 novembre (sans minorant de la fréquence des bonnes idées, rien ne garantit de les
trouver ; les méthodes savantes sont des paris sur une structure). \emph{Pour une paire, trouver la
bonne correspondance} : un problème d'affectation quadratique, difficile même à approcher
\cite{sahni1976} (résumé), dont le paysage a des puits étroits. C'est à ce second niveau que
s'appliquent les méthodes candidates.

\subsection{Les méthodes, une par une}
\begin{itemize}
\item \textbf{Partir des anomalies} : conditionner sur un problème réduit l'espace ; pari de
structure, testable par les prédictions P1 et P3 de la question~7.
\item \textbf{Du grossier au fin par les modes propres} : la matrice $C$ rend les contraintes
linéaires (section~\ref{sec:spectral}) ; la montée spectrale raffine $C$ mais exige une
initialisation et échoue quand elle est mauvaise \cite[p.~11, figure~25, d'après la question~5]{melzi2019}. Les modes grossiers
codent la surface : le grossier lisse le paysage sans rien garantir.
\item \textbf{Relâcher, puis arrondir} : la distance de Gromov et Wasserstein \cite{memoli2011},
régularisée par l'entropie, est non convexe, ne trouve qu'un point stationnaire, à coût
$O(N_1^2N_2+N_2^2N_1)$ au lieu de $O(N_1^2N_2^2)$ \cite[p.~3-4]{peyre2016} ; l'arrondi n'a pas de
garantie. Relâcher, recuire et arrondir sont réunis depuis 1996 dans l'affectation graduée
\cite{goldrangarajan1996} (résumé).
\item \textbf{Recuit} : convergence si et seulement si le refroidissement est logarithmique, de
constante au moins la profondeur du plus profond minimum local non global \cite{hajek1988} (résumé) :
temps exponentiel ; pics très étroits et très larges de masse comparable donnent un mélange lent
\cite[p.~780]{woodard2009}.
\item \textbf{Chercheurs nombreux et divers} : des proposeurs au corpus commun sont corrélés ; générer
séparément, confronter ensuite (section~\ref{sec:cinetique}).
\end{itemize}

La proposition~\ref{prop:lipschitz} précise ce que coûte l'étroitesse locale à une recherche par
transpositions : si le puits est localement large, une montée par transpositions suffit près de
$\varphi$ ; s'il est étroit, il faut des sauts (relâchement, recuit, départs multiples), et la
proposition~\ref{prop:hasard} donne le coût d'un départ au hasard.

\subsection{Amendements proposés (propositions, rien n'est modifié)}

\begin{table}[ht]
\centering\small
\caption{Amendements du Plan Turing proposés par la question~5, section~8 \cite{q5} ; l'erratum de
la question~5 n'en change aucun \cite{erratumq5}.}
\label{tab:amendements}
\begin{tabularx}{\textwidth}{@{}lXl@{}}
\toprule
 & Quoi & Nature \\
\midrule
A1 & une section \enquote{comment chercher} à deux niveaux, qui renvoie aux décisions du 22 novembre et rend obligatoire la base sans modèle & coût nul \\
A2 & un banc de correspondances plantées : $\varphi$ isolée près d'un leurre large de même masse ; relation contre chaîne de deux ; leurre de popularité (modes grossiers communs, structure fine différente) ; bras : énumération exacte sur petits voisinages, criblage de Weisfeiler et Lehman, transport entropique et arrondi, spectral grossier-fin, recuit, proposeurs divers ; mesure : $\varphi$ retrouvées et coût par $\varphi$ & proposer avec recommandation \\
A3 & dans l'épreuve \enquote{le pont suit-il la carte ?}, avant scellement : une spline de distance dans le modèle des naissances, et la permutation des signes à côté de celle des types & coût nul \\
A4 & écrire que \enquote{les cinq signaux localisent les objets nouveaux} est l'énoncé d'un seul modèle, réfuté sous forme lexicale ; toute version relationnelle passe d'abord par un monde de signaux plantés & coût nul \\
A5 & errata de provenance (transition de 2009, jumeaux de Bikard à vérifier, Galison non \enquote{démontré}) & coût nul \\
A6 & les chiffres du bilan de Carbon sont des illustrations, jamais opposables & coût nul \\
\bottomrule
\end{tabularx}
\end{table}

Le résultat négatif du dépôt pèse sur A4 : la rétrodiction des foyers de tension par extraction
lexicale est réfutée (502 articles, onze sous-domaines, aire sous la courbe 0,1667, intervalle à
95~\% de 0 à 0,4667, signal dominé par le registre disciplinaire) ; il réfute un comptage de mots, non
une mesure relationnelle \cite[question~8, point (1)]{reformq8}. Ce qui ne doit pas changer : la file
critique de l'entretien du 1\textsuperscript{er} décembre d'abord ; aucune dépense au premier pas ;
la composition des proposeurs jugée sur l'oracle du monde planté et le coût, jamais sur un juge
\cite[section~8]{q5}.

\begin{statuts}[Section 9 : statuts]
\demontre Rien de neuf (renvois aux propositions~\ref{prop:lipschitz} et \ref{prop:hasard}).
\mesure Rien.
\conjecture Qu'une méthode batte la base sans modèle sur le banc A2, à coût égal ; la question~5 ne
l'affirme pas et n'en fait pas une garantie.
\end{statuts}

%======================================================================
\section{Ce qui est neuf ou non}
\label{sec:neuf}

Règle R8 : une zone qui semble vierge est une question ; la limite d'un article n'est pas un trou du
champ. Rien n'est présenté comme neuf. Les recherches d'antériorité sont celles des pièces sources
(en général deux recherches par question), non refaites ici, sauf pour les références ajoutées par
la note, vérifiées en ligne (notices).

\begin{longtable}{@{}p{0.30\textwidth}p{0.36\textwidth}p{0.28\textwidth}@{}}
\caption{Occupé, ou question ouverte.}\label{tab:neuf}\\
\toprule
Élément & Occupé par & Reste une question \\
\midrule
\endfirsthead
\toprule
Élément & Occupé par & Reste une question \\
\midrule
\endhead
Serrure et clé, systématicité & Gentner 1983 \cite{gentner1983} ; moteur de mise en correspondance \cite{falkenhainer1989} ; équivalence régulière (à vérifier) & l'interaction $S\times Z$ sur la fécondité \\
Cascades A--B, B--C & Swanson 1986 \cite{swanson1986} ; inférence sur prédications \cite{cohen2012} & -- \\
Portée, unification & Friedman 1974 \cite{friedman1974}, Kitcher 1981 \cite{kitcher1981}, Myrvold 2003 \cite{myrvold2003} ; progrès de compression \cite{schmidhuber2010} & leur réunion validée sur monde planté \\
Applications fonctionnelles & Ovsjanikov et collègues 2012 \cite{ovsjanikov2012} ; montée spectrale \cite{melzi2019} & leur emploi entre graphes de domaines typés \\
Appariement relâché & Mémoli 2011 \cite{memoli2011} ; Peyré et collègues 2016 \cite{peyre2016} ; affectation graduée 1996 \cite{goldrangarajan1996} & -- \\
Rupture & Funk et Owen-Smith 2017 \cite{funk2017} ; Wu et collègues 2019 \cite{wu2019} ; Park et collègues 2023 \cite{park2023} & -- \\
Liens manquants & Clauset, Moore et Newman 2008 \cite{clauset2008} & la prédiction d'objets par le cas (C4) \\
Équilibre structurel & Cartwright et Harary 1956 \cite{cartwright1956} ; paysage d'énergie \cite{marvel2009} & lien à la plasticité (question~8) \\
Signaux précurseurs & Scheffer et collègues 2009 \cite{scheffer2009} & sur des graphes de domaines datés (question~8) \\
Poids des problèmes & Laudan 1977 (par Moser \cite{moser1991}) ; Kuhn \cite{kuhn1962} ; questions de Wiśniewski \cite{wisniewski2001} & mesure datée, couplage $\delta P$, monde planté \\
Cinétique, autocatalyse & Arrhenius \cite{arrhenius1889} ; ensembles autocatalytiques \cite{hordijk2004} ; urnes à déclenchement \cite{tria2014} & tracé d'Arrhenius sur des naissances d'idées \\
Transition des coauteurs & Bettencourt, Kaiser et Kaur 2009 \cite{bettencourt2009} & percolation prédite d'une fidélité mesurée \\
Perte des groupes & Diehl et Stroebe 1987 \cite{diehlstroebe1987} & chez des agents parallèles, à coût égal \\
Le pont suit la carte (cortège) & -- (deux recherches de Fable, rien trouvé) & une question, non un trou \\
Co-accélération (résonance) & -- (non cherché ici) & une question ; recherche à faire \\
Clair de lune monstrueux & Conway et Norton 1979 \cite{conwaynorton1979} ; Frenkel, Lepowsky et Meurman 1988 \cite{flm1988} ; Borcherds 1992 \cite{borcherds1992} & question~10, en attente \\
\bottomrule
\end{longtable}

Sur le clair de lune monstrueux, la note vérifie seulement l'arithmétique du complément de Lazar :
$196\,883+1=196\,884$ et $21\,296\,876+196\,883+1=21\,493\,760$ ; le sens des nombres (coefficients de
la fonction modulaire $j$, dimensions de représentations du groupe Monstre) est rapporté, à vérifier
sur la source \cite{conwaynorton1979}.

%======================================================================
\section{Ce que la note ne permet pas de conclure}
\label{sec:limites}

\begin{itemize}
\item Qu'aucune grandeur du modèle prédise de vraies découvertes : aucune épreuve n'a tourné sur
des données réelles ; les seuls chiffres mesurés viennent d'un jouet et de simulations sous
paramètres posés.
\item Que la fusion batte les modèles de Fable et de Carbon ; que les seuils proposés aient la
puissance voulue.
\item Que la \enquote{dimension cachée} existe : l'analogie des corrélations canoniques est une
analogie.
\item Que la profondeur estimée puisse servir de récompense hors d'un monde planté.
\item Quoi que ce soit sur les questions 6, 8 et 10, au-delà du cadre.
\end{itemize}

%======================================================================
\appendix
\section{Sources du dossier et empreintes}
\label{annexe:sources}

\begin{small}
\begin{longtable}{@{}p{0.70\textwidth}l@{}}
\toprule
Pièce (dossier \chemin{docs/notes/}) & empreinte (début) \\
\midrule
\endhead
\chemin{idee-interessante-v1/reponse-fable-max-v1.md} & \texttt{8ec4f9e0b2f8} \\
\chemin{idee-interessante-v1/reponses-google-v1/reponse-carbon-v1.md} & \texttt{a9532fbb60c3} \\
\chemin{idee-interessante-v1/reponses-google-v1/reponse-carbon-homomorphisme-v1.md} & \texttt{3f92cede1557} \\
\chemin{idee-interessante-v1/reponses-google-v1/reponse-carbon-meta-espace-v1.md} & \texttt{01f859c08cb4} \\
\multicolumn{2}{@{}l}{Dans \chemin{questions-directeur-2026-11-24-v1/} :}\\
\chemin{reponse-q1-quatre-etapes-oracles-organisme-v1.md} & \texttt{2cc296327b88} \\
\chemin{reponse-q2-profondeur-surprise-retiree-v1.md} & \texttt{7f4cf73f08f8} \\
\chemin{reponse-q3-cinetique-collisions-idees-v1.md} & \texttt{af126ebffbb9} \\
\chemin{reponse-q3-cinetique-collisions-idees-v1-erratum-1.md} & \texttt{6018b27395e6} \\
\chemin{reponse-q4-emergence-transition-et-groupes-v1.md} & \texttt{249c611c35c9} \\
\chemin{reponse-q4-emergence-transition-et-groupes-v1-erratum-1.md} & \texttt{b359bf154fc3} \\
\chemin{reponse-q5-carbon-fable-et-plan-turing-v1.md} & \texttt{0bdc9810873f} \\
\chemin{reponse-q5-carbon-fable-et-plan-turing-v1-erratum-1.md} & \texttt{7bcaec6f6286} \\
\chemin{reponse-q7-couche-problemes-hypergraphe-v1.md} & \texttt{937e3e693371} \\
\chemin{reponse-q9-expose-graphe-meta-mathematique-v1.md} & \texttt{e981c88c47f4} \\
\chemin{reponse-q9-expose-graphe-meta-mathematique-v1-erratum-1.md} & \texttt{c4a93d16ad13} \\
\chemin{erratum-notations-kappa-et-hypergraphe-v1.md} & \texttt{72afe6930d58} \\
\chemin{errata-mineurs-q1-q2-q7-v1.md} & \texttt{fcac74f59f3d} \\
\chemin{contre-lecture-reponses-q1-q5-q7-q9-v1.md} & \texttt{9cefaa1fab6d} \\
\chemin{complement-note-latex-spectral-puits-fecondite-v1.md} & \texttt{b2dc5ebef071} \\
\chemin{complement-questions-8-9-et-corrections-note-v1.md} & \texttt{58d2e517b07d} \\
\chemin{complement-questions-10-11-et-apport-q5-v1.md} & \texttt{43716a8a79bb} \\
\chemin{demandes-reformulees-quatre-points-q8-q11-v1.md} & \texttt{632c7af28673} \\
\chemin{exemple-graphe-meta-mathematique-v2-sortie.txt} & \texttt{a9d506544bcd} \\
\chemin{simulation-cinetique-barriere-v1-sortie.txt} & \texttt{500e94df2986} \\
\chemin{simulation-cinetique-barriere-v2-sortie.txt} & \texttt{dc98b08f5460} \\
\chemin{simulation-emergence-organisme-v2-sortie.txt} & \texttt{a31ee52ca0c8} \\
\bottomrule
\end{longtable}
\end{small}

Fichiers lus hors du dépôt par la note (bloc-notes de la session, droits des éditeurs) : Wu, Wang et
Evans, arXiv 1709.02445v1 (\texttt{51a48bd3909a}), p.~7 et 18 ; Ovsjanikov et collègues 2012
(\texttt{7d7fb34a3e61}), p.~1 à 6.

%======================================================================
\begin{thebibliography}{99}\small
\setlength{\itemsep}{1pt}

\bibitem{fable} Fable 5.2 (23 novembre 2026). Réponse : modèle testable de l'idée intéressante et
de sa naissance, v1. \chemin{docs/notes/idee-interessante-v1/reponse-fable-max-v1.md}
(\texttt{8ec4f9e0…}).
\bibitem{carbon} Carbon. Réponse sur le modèle de l'idée intéressante, v1.
\chemin{docs/notes/idee-interessante-v1/reponses-google-v1/reponse-carbon-v1.md} (\texttt{a9532fbb…}) ;
références de Carbon non vérifiées ici.
\bibitem{carbonhomo} Carbon. Réponse sur l'homomorphisme, v1. Même dossier,
\chemin{reponse-carbon-homomorphisme-v1.md} (\texttt{3f92cede…}).
\bibitem{carbonmeta} Carbon. Réponse sur le méta-espace mathématique, v1. Même dossier,
\chemin{reponse-carbon-meta-espace-v1.md} (\texttt{01f859c0…}).
\bibitem{q2} Réponse à la question 2 (profondeur comme surprise retirée), v1 (\texttt{7f4cf73f…}).
\bibitem{q3} Réponse à la question 3 (cinétique des collisions d'idées), v1 (\texttt{af126ebf…}).
\bibitem{erratumq3} Erratum 1 de la question 3 (\texttt{6018b273…}).
\bibitem{q4} Réponse à la question 4 (émergence), v1 (\texttt{249c611c…}).
\bibitem{erratumq4} Erratum 1 de la question 4 (\texttt{b359bf15…}).
\bibitem{q5} Réponse à la question 5 (Carbon contre Fable ; Plan Turing), v1 (\texttt{0bdc9810…}).
\bibitem{erratumq5} Erratum 1 de la question 5 (\texttt{7bcaec6f…}).
\bibitem{q7} Réponse à la question 7 (couche des problèmes, hypergraphe), v1 (\texttt{937e3e69…}).
\bibitem{q9} Exposé du graphe méta-mathématique (question 9), v1 (\texttt{e981c88c…}).
\bibitem{erratumq9} Erratum 1 de la question 9 (\texttt{c4a93d16…}).
\bibitem{sortieq9v2} Sortie de \chemin{exemple-graphe-meta-mathematique-v2.py} (\texttt{a9d50654…}).
\bibitem{erratumkappa} Erratum transversal : une seule notation pour $\kappa$, hypergraphe fusionné
(\texttt{72afe693…}).
\bibitem{erratamineurs} Errata mineurs des questions 1, 2 et 7 (\texttt{fcac74f5…}).
\bibitem{contrelecture} Contre-lecture des réponses aux questions 1 à 5, 7 et 9 (\texttt{9cefaa1f…}).
\bibitem{complementspectral} Complément : coordonnées spectrales, puits profond et étroit,
fécondité (\texttt{b2dc5ebe…}).
\bibitem{complement89} Complément : questions 8 et 9, corrections pour la note (\texttt{58d2e517…}).
\bibitem{reformq8} Demandes reformulées en quatre points, questions 8 à 11 (\texttt{632c7af2…}).

\bibitem{abi2021} Abi, B. et al. (2021). Measurement of the positive muon anomalous magnetic moment
to 0.46 ppm. \emph{Physical Review Letters} 126, 141801. \url{https://doi.org/10.1103/PhysRevLett.126.141801}
(résumé, lu par la question 7).
\bibitem{aliberti2025} Aliberti, R. et al. (2025). The anomalous magnetic moment of the muon in the
Standard Model: an update. \emph{Physics Reports} 1143, 1-158.
\url{https://doi.org/10.1016/j.physrep.2025.08.002} (résumé, lu par la question 7).
\bibitem{arrhenius1889} Arrhenius, S. (1889). Über die Reaktionsgeschwindigkeit bei der Inversion von
Rohrzucker durch Säuren. \emph{Zeitschrift für Physikalische Chemie} 4, 226-248.
\url{https://doi.org/10.1515/zpch-1889-0416} (notice).
\bibitem{bettencourt2009} Bettencourt, L. M. A., Kaiser, D. I. et Kaur, J. (2009). Scientific discovery
and topological transitions in collaboration networks. \emph{Journal of Informetrics} 3(3), 210-221.
\url{https://doi.org/10.1016/j.joi.2009.03.001} (manuscrit d'auteur lu par la question 4, p.~5 et 14).
\bibitem{borcherds1992} Borcherds, R. E. (1992). Monstrous moonshine and monstrous Lie superalgebras.
\emph{Inventiones Mathematicae} 109(1), 405-444. \url{https://doi.org/10.1007/BF01232032} (notice
Crossref vérifiée le 24 novembre 2026).
\bibitem{bornmann2020} Bornmann, L., Devarakonda, S., Tekles, A. et Chacko, G. (2020). Are disruption
index indicators convergently valid? \emph{Quantitative Science Studies} 1(3), 1242-1259.
\url{https://doi.org/10.1162/qss_a_00068} (résumé, lu par la question 2).
\bibitem{boudreau2016} Boudreau, K. J. et al. (2016). Looking across and looking beyond the knowledge
frontier. \emph{Management Science} 62, 2765-2783. \url{https://doi.org/10.1287/mnsc.2015.2285}
(cité par Fable).
\bibitem{cartwright1956} Cartwright, D. et Harary, F. (1956). Structural balance: a generalization of
Heider's theory. \emph{Psychological Review} 63(5), 277-293. \url{https://doi.org/10.1037/h0046049}
(notice Crossref vérifiée).
\bibitem{chung1997} Chung, F. R. K. (1997). \emph{Spectral Graph Theory}. CBMS Regional Conference
Series in Mathematics 92, American Mathematical Society. \url{https://doi.org/10.1090/cbms/092}
(notice Crossref vérifiée, datée de décembre 1996 ; définitions rappelées, pages non citées).
\bibitem{chung2005} Chung, F. (2005). Laplacians and the Cheeger inequality for directed graphs.
\emph{Annals of Combinatorics} 9(1), 1-19. \url{https://doi.org/10.1007/s00026-005-0237-z} (notice
Crossref vérifiée ; contenu rappelé, à vérifier sur le texte).
\bibitem{clauset2008} Clauset, A., Moore, C. et Newman, M. E. J. (2008). Hierarchical structure and the
prediction of missing links in networks. \emph{Nature} 453, 98-101.
\url{https://doi.org/10.1038/nature06830} (notice Crossref vérifiée).
\bibitem{cohen2012} Cohen, T. et al. (2012). Discovering discovery patterns with predication-based
semantic indexing. \emph{Journal of Biomedical Informatics} 45, 1049-1065.
\url{https://doi.org/10.1016/j.jbi.2012.07.003} (cité par Fable).
\bibitem{collins1949} Collins, F. C. et Kimball, G. E. (1949). Diffusion-controlled reaction rates.
\emph{Journal of Colloid Science} 4, 425-437. \url{https://doi.org/10.1016/0095-8522(49)90023-9}
(notice ; forme rapportée, page à vérifier).
\bibitem{conwaynorton1979} Conway, J. H. et Norton, S. P. (1979). Monstrous moonshine. \emph{Bulletin of
the London Mathematical Society} 11(3), 308-339. \url{https://doi.org/10.1112/blms/11.3.308} (notice
Crossref vérifiée ; contenu à vérifier).
\bibitem{davis1971} Davis, M. S. (1971). That's interesting! \emph{Philosophy of the Social Sciences} 1,
309-344. \url{https://doi.org/10.1177/004839317100100211} (cité par Fable).
\bibitem{daviskahan1970} Davis, C. et Kahan, W. M. (1970). The rotation of eigenvectors by a
perturbation. III. \emph{SIAM Journal on Numerical Analysis} 7(1), 1-46.
\url{https://doi.org/10.1137/0707001} (notice Crossref vérifiée ; énoncé rappelé, page à vérifier).
\bibitem{diehlstroebe1987} Diehl, M. et Stroebe, W. (1987). Productivity loss in brainstorming groups:
toward the solution of a riddle. \emph{Journal of Personality and Social Psychology} 53(3), 497-509.
\url{https://doi.org/10.1037/0022-3514.53.3.497} (PDF lu par la question 4, p.~497).
\bibitem{evanspolanyi1938} Evans, M. G. et Polanyi, M. (1938). Inertia and driving force of chemical
reactions. \emph{Transactions of the Faraday Society} 34, 11. \url{https://doi.org/10.1039/tf9383400011}
(notice ; relation rapportée, page à vérifier).
\bibitem{falkenhainer1989} Falkenhainer, B., Forbus, K. D. et Gentner, D. (1989). The structure-mapping
engine. \emph{Artificial Intelligence} 41, 1-63. \url{https://doi.org/10.1016/0004-3702(89)90077-5}
(cité par Fable).
\bibitem{fiedler1973} Fiedler, M. (1973). Algebraic connectivity of graphs. \emph{Czechoslovak
Mathematical Journal} 23(2), 298-305. \url{https://doi.org/10.21136/CMJ.1973.101168} (notice Crossref
vérifiée).
\bibitem{flm1988} Frenkel, I., Lepowsky, J. et Meurman, A. (1988). \emph{Vertex Operator Algebras and
the Monster}. Pure and Applied Mathematics 134, Academic Press. Notice Crossref de la préface :
\url{https://doi.org/10.1016/S0079-8169(08)62324-7} (vérifiée ; contenu à vérifier).
\bibitem{friedman1974} Friedman, M. (1974). Explanation and scientific understanding. \emph{The Journal
of Philosophy} 71(1), 5-19. \url{https://doi.org/10.2307/2024924} (texte lu par la question 2, p.~15).
\bibitem{funk2017} Funk, R. J. et Owen-Smith, J. (2017). A dynamic network measure of technological
change. \emph{Management Science} 63(3), 791-817. \url{https://doi.org/10.1287/mnsc.2015.2366} (notice
Crossref vérifiée, résumé lu par la question 2 ; définition exacte de l'indice à vérifier).
\bibitem{gallo1993} Gallo, G., Longo, G., Pallottino, S. et Nguyen, S. (1993). Directed hypergraphs and
applications. \emph{Discrete Applied Mathematics} 42(2-3), 177-201.
\url{https://doi.org/10.1016/0166-218X(93)90045-P} (notice Crossref vérifiée ; résumé).
\bibitem{gentner1983} Gentner, D. (1983). Structure-mapping: a theoretical framework for analogy.
\emph{Cognitive Science} 7, 155-170. \url{https://doi.org/10.1207/s15516709cog0702_3} (cité par Fable).
\bibitem{goldrangarajan1996} Gold, S. et Rangarajan, A. (1996). A graduated assignment algorithm for
graph matching. \emph{IEEE Transactions on Pattern Analysis and Machine Intelligence} 18(4), 377-388.
\url{https://doi.org/10.1109/34.491619} (résumé, lu par la question 5).
\bibitem{gukrenn2026} Gu, X. et Krenn, M. (2026). Generating interesting scientific ideas using
knowledge graphs and LLMs: evaluations with 100 research group leaders. \emph{Proceedings of the
National Academy of Sciences} 123(46), e2611349123. \url{https://doi.org/10.1073/pnas.2611349123} ;
arXiv 2405.17044v4 (PDF lu par la question 2, p.~8 ; chiffres confirmés par la contre-lecture).
\bibitem{hajek1988} Hajek, B. (1988). Cooling schedules for optimal annealing. \emph{Mathematics of
Operations Research} 13(2), 311-329. \url{https://doi.org/10.1287/moor.13.2.311} (résumé, lu par la
question 5).
\bibitem{harris1963} Harris, T. E. (1963). \emph{The Theory of Branching Processes}. Springer.
\url{https://doi.org/10.1007/978-3-642-51866-9} (notice Crossref vérifiée ; théorème d'extinction
rappelé, page à vérifier).
\bibitem{hinton2015} Hinton, G., Vinyals, O. et Dean, J. (2015). Distilling the knowledge in a neural
network. arXiv 1503.02531v1 (PDF lu par la question 3, éq.~(1), p.~2-3).
\bibitem{hordijk2004} Hordijk, W. et Steel, M. (2004). Detecting autocatalytic, self-sustaining sets in
chemical reaction systems. \emph{Journal of Theoretical Biology} 227, 451-461.
\url{https://doi.org/10.1016/j.jtbi.2003.11.020} (résumé, lu par la question 3).
\bibitem{hotelling1936} Hotelling, H. (1936). Relations between two sets of variates. \emph{Biometrika}
28(3/4), 321-377. \url{https://doi.org/10.2307/2333955} (notice Crossref vérifiée ; définition
rappelée, texte non lu).
\bibitem{kitcher1981} Kitcher, P. (1981). Explanatory unification. \emph{Philosophy of Science} 48(4),
507-531. \url{https://doi.org/10.1086/289019} (texte lu par la question 2, p.~514 et 526-529).
\bibitem{kuhn1962} Kuhn, T. S. (1962 ; 2\textsuperscript{e} éd. 1970). \emph{The Structure of Scientific
Revolutions}. University of Chicago Press. \url{https://doi.org/10.7208/chicago/9780226458106.001.0001}
(extraits lus par la question 7, p.~52-53 et 81-82).
\bibitem{kunegis2010} Kunegis, J., Schmidt, S., Lommatzsch, A., Lerner, J., De Luca, E. W. et Albayrak,
S. (2010). Spectral analysis of signed graphs for clustering, prediction and visualization.
\emph{Proceedings of the 2010 SIAM International Conference on Data Mining}, 559-570.
\url{https://doi.org/10.1137/1.9781611972801.49} (notice Crossref vérifiée).
\bibitem{leibel2024} Leibel, C. et Bornmann, L. (2024). What do we know about the disruption index in
scientometrics? \emph{Scientometrics} 129, 601-639. \url{https://doi.org/10.1007/s11192-023-04873-5}
(résumé, lu par la question 2).
\bibitem{marvel2009} Marvel, S. A., Strogatz, S. H. et Kleinberg, J. M. (2009). Energy landscape of
social balance. \emph{Physical Review Letters} 103, 198701.
\url{https://doi.org/10.1103/PhysRevLett.103.198701} (notice Crossref vérifiée).
\bibitem{melzi2019} Melzi, S., Ren, J., Rodolà, E., Sharma, A., Wonka, P. et Ovsjanikov, M. (2019).
ZoomOut: spectral upsampling for efficient shape correspondence. \emph{ACM Transactions on Graphics}
38(6). \url{https://doi.org/10.1145/3355089.3356524} ; arXiv 1904.07865 (PDF lu par la question 5).
\bibitem{memoli2011} Mémoli, F. (2011). Gromov-Wasserstein distances and the metric approach to object
matching. \emph{Foundations of Computational Mathematics} 11(4), 417-487.
\url{https://doi.org/10.1007/s10208-011-9093-5} (notice).
\bibitem{molloyreed1995} Molloy, M. et Reed, B. (1995). A critical point for random graphs with a given
degree sequence. \emph{Random Structures and Algorithms} 6(2-3), 161-180.
\url{https://doi.org/10.1002/rsa.3240060204} (résumé, lu par la question 4).
\bibitem{moser1991} Moser, P. K. (1991). Scientific rationality and progress: a paradox in Laudan's
account. \emph{Theoria et Historia Scientiarum} 1, 75-81. \url{https://doi.org/10.12775/ths.1991.006}
(PDF lu par la question 7, p.~76-78 ; Laudan 1977, livre non vu).
\bibitem{mu2018} Mu, J., Bhat, S. et Viswanath, P. (2018). All-but-the-top: simple and effective
postprocessing for word representations. \emph{ICLR 2018}. arXiv 1702.01417v2 (p.~4, lue par Fable).
\bibitem{myrvold2003} Myrvold, W. C. (2003). A Bayesian account of the virtue of unification.
\emph{Philosophy of Science} 70(2), 399-423. \url{https://doi.org/10.1086/375475} (équations lues dans
le manuscrit d'auteur par la question 2 ; pages publiées à vérifier).
\bibitem{newman2001pre} Newman, M. E. J., Strogatz, S. H. et Watts, D. J. (2001). Random graphs with
arbitrary degree distributions and their applications. \emph{Physical Review E} 64, 026118.
\url{https://doi.org/10.1103/PhysRevE.64.026118} (résumé, lu par la question 4).
\bibitem{ovsjanikov2012} Ovsjanikov, M., Ben-Chen, M., Solomon, J., Butscher, A. et Guibas, L. (2012).
Functional maps: a flexible representation of maps between shapes. \emph{ACM Transactions on
Graphics} 31(4), 30. \url{https://doi.org/10.1145/2185520.2185526} (PDF lu, p.~1-6 ; empreinte
\texttt{7d7fb34a…}).
\bibitem{park2023} Park, M., Leahey, E. et Funk, R. J. (2023). Papers and patents are becoming less
disruptive over time. \emph{Nature} 613, 138-144. \url{https://doi.org/10.1038/s41586-022-05543-x}
(notice Crossref vérifiée ; texte non lu).
\bibitem{peyre2016} Peyré, G., Cuturi, M. et Solomon, J. (2016). Gromov-Wasserstein averaging of kernel
and distance matrices. \emph{ICML 2016}, PMLR 48, 2664-2672 (PDF lu par la question 5, p.~2-4).
\bibitem{pubtator2024} Wei, C.-H. et al. (2024). PubTator 3.0: an AI-powered literature resource for
unlocking biomedical knowledge. \emph{Nucleic Acids Research} 52, W540-W546.
\url{https://doi.org/10.1093/nar/gkae235} (texte lu par la question 7).
\bibitem{rissanen1978} Rissanen, J. (1978). Modeling by shortest data description. \emph{Automatica}
14(5), 465-471. \url{https://doi.org/10.1016/0005-1098(78)90005-5} (notice).
\bibitem{rosemblat2019} Rosemblat, G., Fiszman, M., Shin, D. et Kilicoglu, H. (2019). Towards a
characterization of apparent contradictions in the biomedical literature using context analysis.
\emph{Journal of Biomedical Informatics} 98, 103275. \url{https://doi.org/10.1016/j.jbi.2019.103275}
(résumé, lu par la question 7).
\bibitem{sahni1976} Sahni, S. et Gonzalez, T. (1976). P-complete approximation problems. \emph{Journal of
the ACM} 23(3), 555-565. \url{https://doi.org/10.1145/321958.321975} (résumé, lu par la question 5).
\bibitem{scheffer2009} Scheffer, M. et al. (2009). Early-warning signals for critical transitions.
\emph{Nature} 461, 53-59. \url{https://doi.org/10.1038/nature08227} (notice Crossref vérifiée).
\bibitem{schmidhuber2010} Schmidhuber, J. (2010). Formal theory of creativity, fun, and intrinsic
motivation (1990-2010). \emph{IEEE Transactions on Autonomous Mental Development} 2(3), 230-247.
\url{https://doi.org/10.1109/TAMD.2010.2056368} (prépublication lue par la question 2, p.~3 et 5).
\bibitem{shievans2023} Shi, F. et Evans, J. (2023). Surprising combinations of research contents and
contexts are related to impact and emerge with scientific outsiders from distant disciplines.
\emph{Nature Communications} 14, 1641. \url{https://doi.org/10.1038/s41467-023-36741-4} (résumé).
\bibitem{swanson1986} Swanson, D. R. (1986). Fish oil, Raynaud's syndrome, and undiscovered public
knowledge. \emph{Perspectives in Biology and Medicine} 30, 7-18.
\url{https://doi.org/10.1353/pbm.1986.0087} (cité par Fable).
\bibitem{tria2014} Tria, F., Loreto, V., Servedio, V. D. P. et Strogatz, S. H. (2014). The dynamics of
correlated novelties. \emph{Scientific Reports} 4, 5890. \url{https://doi.org/10.1038/srep05890} (PDF lu
par la question 3, p.~3, légende de la figure~1).
\bibitem{wisniewski2001} Wiśniewski, A. (2001). Questions and inferences. \emph{Logique et Analyse}
173-175, 5-43 (PDF lu par la question 7, définition 8, p.~17).
\bibitem{woodard2009} Woodard, D. B., Schmidler, S. C. et Huber, M. (2009). Sufficient conditions for
torpid mixing of parallel and simulated tempering. \emph{Electronic Journal of Probability} 14,
780-804. \url{https://doi.org/10.1214/EJP.v14-638} (p.~780 vérifiée sur PDF par une pièce antérieure,
d'après la question 5).
\bibitem{wu2019} Wu, L., Wang, D. et Evans, J. A. (2019). Large teams develop and small teams disrupt
science and technology. \emph{Nature} 566, 378-382. \url{https://doi.org/10.1038/s41586-019-0941-9} ;
prépublication arXiv 1709.02445v1 (PDF lu par la note, p.~7, légende de la figure~1, et p.~18,
figure~S3 ; empreinte \texttt{51a48bd3…} ; version publiée non lue).

\end{thebibliography}

\end{document}
